% Where the loops meet: each way loops can meet, three times over: the circuit you build, the loop it is, the sound you hear.
% The heavy stroke is the path that comes back: a feedback wire in the circuit, a circle on its own.
% Geometry: gen-loops.py (torus-mesh.tex, torus-path.tex, stretch-rings.tex). Data: gen-where-the-loops-meet.py.
% Rows 1-3 are computed (ink); row 4's curve is measured by ramp.mjs on the golden pair (measure red).
% Compile: lualatex where-the-loops-meet.tex ; SVG via PyMuPDF (dvisvgm's PDF route is broken on this Mac).
\documentclass[border=0pt]{standalone}
\usepackage{fontspec}
\setmainfont{JetBrains Mono}
\setmonofont{JetBrains Mono}
\usepackage{tikz}
\usepackage{pgfplots}
\pgfplotsset{compat=1.16}
\usetikzlibrary{arrows.meta, calc, backgrounds, decorations.pathmorphing, decorations.markings, positioning}
% Drafting Table, paper ground (_ops/drafting-table/design-system/colors_and_type.css)
\definecolor{paper}{HTML}{F0EEE5}
\definecolor{ink}{HTML}{1A1916}
\definecolor{ink2}{HTML}{3A3A35}
\definecolor{ink3}{HTML}{6A6862}
\definecolor{slab}{HTML}{D8D4C6}
\definecolor{measure}{HTML}{C8261A}
\tikzset{
  >={Stealth[length=1.8mm]},
  mesh/.style={draw=ink3!45, line width=0.25pt},
  scan/.style={draw=ink, line width=0.6pt},
  scanback/.style={draw=ink3!70, line width=0.4pt, densely dashed},
  ring/.style={draw=ink, line width=0.9pt},
  loop/.style={draw=ink, line width=1.5pt},
  back/.style={draw=ink, line width=1.5pt, ->},
  wire/.style={draw=ink, line width=0.6pt},
  blk/.style={rectangle, draw=ink, fill=slab, line width=0.7pt, minimum height=7mm, minimum width=12mm, inner sep=2pt, font=\scriptsize, text=ink},
  adder/.style={circle, draw=ink, fill=paper, line width=0.7pt, minimum size=4.6mm, inner sep=0pt, font=\scriptsize},
  tap/.style={circle, fill=ink, minimum size=3pt, inner sep=0pt},
  lab/.style={font=\scriptsize, text=ink2},
  rowhead/.style={font=\bfseries\small, anchor=west, text=ink},
  rowsub/.style={font=\scriptsize, anchor=west, text=ink2},
  colhead/.style={font=\bfseries\scriptsize, anchor=west, text=ink},
  spin/.style={postaction={decorate}, decoration={markings, mark=at position 0.25 with {\arrow{>}}}},
}
\pgfplotsset{
  /pgf/number format/assume math mode=true,
  snd/.style={width=6.6cm, height=3.0cm, axis lines=left, tick align=outside, ymin=0, ytick=\empty,
    axis line style={ink, line width=0.5pt}, tick style={ink}, clip=false,
    xlabel style={font=\scriptsize, text=ink2, yshift=1mm}, ylabel style={font=\scriptsize, text=ink2},
    tick label style={font=\tiny, text=ink2}},
}
% a row: frame, heading, and the two column rules
\newcommand{\rowframe}[3]{%
  \draw[ink, line width=0.5pt] (-0.3,#1) rectangle (19.9,#1-5.9);
  \draw[ink, line width=0.3pt] (6.9,#1-0.95) -- (6.9,#1-5.9) (12.4,#1-0.95) -- (12.4,#1-5.9) (-0.3,#1-0.95) -- (19.9,#1-0.95);
  \node[rowhead] at (0,#1-0.32) {#2};
  \node[rowsub] at (0,#1-0.7) {#3};}
\begin{document}
\begin{tikzpicture}

% ---------- column heads ----------
\node[colhead] at (0,1.15) {THE CIRCUIT \quad what you build};
\node[colhead] at (7.2,1.15) {THE LOOP \quad what it is};
\node[colhead] at (12.7,1.15) {THE SOUND \quad what you hear};
\node[font=\scriptsize, anchor=west, text=ink2] at (0,0.55) {Heavy: the path that comes back. In the circuit it is a feedback wire or a phasor that wraps; on its own it is a circle.};

% =============== 01 · READ THROUGH ONE TABLE ===============
\def\y{0}
\rowframe{\y}{01 \enspace READ THROUGH ONE TABLE}{a lattice in frequency: every m·f₁ + n·f₂}
\begin{scope}[shift={(0,\y)}]
  \node[blk] (p1) at (1.2,-2.5) {phasor f₁};
  \node[blk] (p2) at (1.2,-4.85) {phasor f₂};
  \draw[back] (p1.north east) to[out=40, in=140, looseness=2.2] node[lab, above] {wraps} (p1.north west);
  \draw[back] (p2.north east) to[out=40, in=140, looseness=2.2] node[lab, above] {wraps} (p2.north west);
  \node[blk, minimum width=20mm, minimum height=12mm] (w) at (4.3,-3.7) {table W(φ₁, φ₂)};
  \draw[wire,->] (p1.east) -| ([xshift=-4mm]w.north);
  \draw[wire,->] (p2.east) -| ([xshift=-4mm]w.south);
  \draw[wire,->] (w.east) -- (6.65,-3.7) node[lab, above left] {out};
  % the loop: two circles, read together, are a torus
  \draw[loop, spin] (7.75,-2.05) circle (0.42); \node[lab] at (7.75,-2.05) {φ₁};
  \draw[loop, spin] (7.75,-4.2) circle (0.3);  \node[lab] at (7.75,-4.2) {φ₂};
  \node[font=\small, text=ink] at (8.35,-3.15) {×};
  \begin{scope}[shift={(10.45,-3.25)}, scale=0.75]
    % torus mesh, generated by gen-loops.py
\draw[mesh] (2.428,0.042) -- (2.422,0.084) -- (2.411,0.125) -- (2.397,0.166) -- (2.378,0.205) -- (2.356,0.243) -- (2.330,0.280) -- (2.300,0.315) -- (2.267,0.348) -- (2.231,0.379) -- (2.192,0.407) -- (2.150,0.434) -- (2.105,0.457) -- (2.059,0.477) -- (2.010,0.495) -- (1.960,0.510) -- (1.909,0.521) -- (1.856,0.529) -- (1.803,0.534) -- (1.750,0.536) -- (1.697,0.534) -- (1.644,0.529) -- (1.591,0.521) -- (1.540,0.510) -- (1.490,0.495) -- (1.441,0.477) -- (1.395,0.457) -- (1.350,0.434) -- (1.308,0.407) -- (1.269,0.379) -- (1.233,0.348) -- (1.200,0.315) -- (1.170,0.280) -- (1.144,0.243) -- (1.122,0.205) -- (1.103,0.166) -- (1.089,0.125) -- (1.078,0.084) -- (1.072,0.042) -- (1.070,0.000) -- (1.072,-0.042);
\draw[mesh] (2.297,0.584) -- (2.276,0.619) -- (2.250,0.651) -- (2.222,0.681) -- (2.190,0.709) -- (2.155,0.734) -- (2.117,0.757) -- (2.076,0.776) -- (2.034,0.792) -- (1.989,0.805) -- (1.942,0.815) -- (1.893,0.822) -- (1.844,0.825) -- (1.793,0.825) -- (1.742,0.822) -- (1.690,0.815) -- (1.639,0.805) -- (1.588,0.791) -- (1.537,0.775) -- (1.487,0.755) -- (1.439,0.732) -- (1.392,0.707) -- (1.347,0.679) -- (1.304,0.649) -- (1.264,0.616) -- (1.226,0.581) -- (1.191,0.544) -- (1.159,0.506) -- (1.130,0.466) -- (1.105,0.426) -- (1.084,0.384) -- (1.066,0.341) -- (1.052,0.299) -- (1.042,0.256) -- (1.036,0.213) -- (1.034,0.170) -- (1.036,0.129) -- (1.042,0.088) -- (1.052,0.048) -- (1.066,0.010) -- (1.084,-0.026);
\draw[mesh] (1.992,1.023) -- (1.963,1.046) -- (1.932,1.066) -- (1.898,1.082) -- (1.862,1.095) -- (1.823,1.105) -- (1.783,1.111) -- (1.741,1.114) -- (1.698,1.113) -- (1.653,1.109) -- (1.608,1.101) -- (1.562,1.089) -- (1.516,1.075) -- (1.469,1.056) -- (1.423,1.035) -- (1.378,1.011) -- (1.334,0.984) -- (1.290,0.954) -- (1.248,0.921) -- (1.208,0.886) -- (1.169,0.849) -- (1.133,0.810) -- (1.099,0.770) -- (1.068,0.728) -- (1.039,0.684) -- (1.013,0.640) -- (0.991,0.595) -- (0.971,0.550) -- (0.955,0.505) -- (0.943,0.460) -- (0.934,0.416) -- (0.928,0.372) -- (0.927,0.329) -- (0.928,0.288) -- (0.934,0.248) -- (0.943,0.210) -- (0.955,0.174) -- (0.971,0.140) -- (0.991,0.109) -- (1.013,0.080) -- (1.039,0.054);
\draw[mesh] (1.577,1.350) -- (1.550,1.362) -- (1.520,1.369) -- (1.489,1.373) -- (1.456,1.374) -- (1.421,1.370) -- (1.386,1.363) -- (1.350,1.352) -- (1.313,1.337) -- (1.275,1.319) -- (1.237,1.298) -- (1.200,1.273) -- (1.162,1.245) -- (1.125,1.214) -- (1.089,1.180) -- (1.053,1.144) -- (1.019,1.105) -- (0.986,1.064) -- (0.955,1.021) -- (0.925,0.977) -- (0.897,0.931) -- (0.872,0.885) -- (0.848,0.837) -- (0.827,0.789) -- (0.809,0.741) -- (0.793,0.693) -- (0.780,0.646) -- (0.770,0.599) -- (0.763,0.553) -- (0.758,0.509) -- (0.757,0.466) -- (0.758,0.425) -- (0.763,0.386) -- (0.770,0.349) -- (0.780,0.315) -- (0.793,0.283) -- (0.809,0.255) -- (0.827,0.229) -- (0.848,0.207) -- (0.872,0.189) -- (0.897,0.174);
\draw[mesh] (1.096,1.576) -- (1.075,1.580) -- (1.053,1.579) -- (1.029,1.575) -- (1.005,1.567) -- (0.980,1.555) -- (0.954,1.539) -- (0.928,1.519) -- (0.902,1.496) -- (0.875,1.469) -- (0.848,1.439) -- (0.822,1.406) -- (0.796,1.369) -- (0.770,1.331) -- (0.745,1.289) -- (0.721,1.246) -- (0.697,1.201) -- (0.675,1.153) -- (0.654,1.105) -- (0.635,1.056) -- (0.616,1.005) -- (0.600,0.955) -- (0.585,0.904) -- (0.572,0.853) -- (0.561,0.803) -- (0.552,0.754) -- (0.544,0.706) -- (0.539,0.659) -- (0.536,0.614) -- (0.535,0.571) -- (0.536,0.530) -- (0.539,0.491) -- (0.544,0.455) -- (0.552,0.423) -- (0.561,0.393) -- (0.572,0.367) -- (0.585,0.344) -- (0.600,0.325) -- (0.616,0.309) -- (0.635,0.298) -- (0.654,0.290);
\draw[mesh] (0.556,1.712) -- (0.545,1.709) -- (0.533,1.702) -- (0.520,1.691) -- (0.507,1.675) -- (0.494,1.656) -- (0.480,1.633) -- (0.467,1.607) -- (0.453,1.577) -- (0.439,1.543) -- (0.425,1.507) -- (0.412,1.467) -- (0.399,1.425) -- (0.386,1.381) -- (0.373,1.335) -- (0.361,1.286) -- (0.349,1.237) -- (0.339,1.186) -- (0.328,1.134) -- (0.319,1.081) -- (0.311,1.029) -- (0.303,0.976) -- (0.296,0.924) -- (0.290,0.872) -- (0.286,0.822) -- (0.282,0.773) -- (0.279,0.725) -- (0.277,0.680) -- (0.277,0.636) -- (0.277,0.596) -- (0.279,0.557) -- (0.282,0.522) -- (0.286,0.491) -- (0.290,0.462) -- (0.296,0.437) -- (0.303,0.416) -- (0.311,0.399) -- (0.319,0.385) -- (0.328,0.376) -- (0.339,0.371) -- (0.349,0.369);
\draw[mesh] (0.000,1.757) -- (0.000,1.753) -- (0.000,1.745) -- (0.000,1.733) -- (0.000,1.716) -- (0.000,1.696) -- (0.000,1.672) -- (0.000,1.644) -- (0.000,1.613) -- (0.000,1.579) -- (0.000,1.541) -- (0.000,1.501) -- (0.000,1.458) -- (0.000,1.412) -- (0.000,1.365) -- (0.000,1.316) -- (0.000,1.265) -- (0.000,1.213) -- (0.000,1.160) -- (0.000,1.107) -- (0.000,1.054) -- (0.000,1.000) -- (0.000,0.948) -- (0.000,0.896) -- (0.000,0.845) -- (0.000,0.795) -- (0.000,0.748) -- (0.000,0.702) -- (0.000,0.659) -- (0.000,0.618) -- (0.000,0.580) -- (0.000,0.545) -- (0.000,0.514) -- (0.000,0.486) -- (0.000,0.461) -- (0.000,0.440) -- (0.000,0.424) -- (0.000,0.411) -- (0.000,0.402) -- (0.000,0.398) -- (0.000,0.398);
\draw[mesh] (-0.556,1.712) -- (-0.545,1.709) -- (-0.533,1.702) -- (-0.520,1.691) -- (-0.507,1.675) -- (-0.494,1.656) -- (-0.480,1.633) -- (-0.467,1.607) -- (-0.453,1.577) -- (-0.439,1.543) -- (-0.425,1.507) -- (-0.412,1.467) -- (-0.399,1.425) -- (-0.386,1.381) -- (-0.373,1.335) -- (-0.361,1.286) -- (-0.349,1.237) -- (-0.339,1.186) -- (-0.328,1.134) -- (-0.319,1.081) -- (-0.311,1.029) -- (-0.303,0.976) -- (-0.296,0.924) -- (-0.290,0.872) -- (-0.286,0.822) -- (-0.282,0.773) -- (-0.279,0.725) -- (-0.277,0.680) -- (-0.277,0.636) -- (-0.277,0.596) -- (-0.279,0.557) -- (-0.282,0.522) -- (-0.286,0.491) -- (-0.290,0.462) -- (-0.296,0.437) -- (-0.303,0.416) -- (-0.311,0.399) -- (-0.319,0.385) -- (-0.328,0.376) -- (-0.339,0.371) -- (-0.349,0.369);
\draw[mesh] (-1.096,1.576) -- (-1.075,1.580) -- (-1.053,1.579) -- (-1.029,1.575) -- (-1.005,1.567) -- (-0.980,1.555) -- (-0.954,1.539) -- (-0.928,1.519) -- (-0.902,1.496) -- (-0.875,1.469) -- (-0.848,1.439) -- (-0.822,1.406) -- (-0.796,1.369) -- (-0.770,1.331) -- (-0.745,1.289) -- (-0.721,1.246) -- (-0.697,1.201) -- (-0.675,1.153) -- (-0.654,1.105) -- (-0.635,1.056) -- (-0.616,1.005) -- (-0.600,0.955) -- (-0.585,0.904) -- (-0.572,0.853) -- (-0.561,0.803) -- (-0.552,0.754) -- (-0.544,0.706) -- (-0.539,0.659) -- (-0.536,0.614) -- (-0.535,0.571) -- (-0.536,0.530) -- (-0.539,0.491) -- (-0.544,0.455) -- (-0.552,0.423) -- (-0.561,0.393) -- (-0.572,0.367) -- (-0.585,0.344) -- (-0.600,0.325) -- (-0.616,0.309) -- (-0.635,0.298) -- (-0.654,0.290);
\draw[mesh] (-1.577,1.350) -- (-1.550,1.362) -- (-1.520,1.369) -- (-1.489,1.373) -- (-1.456,1.374) -- (-1.421,1.370) -- (-1.386,1.363) -- (-1.350,1.352) -- (-1.313,1.337) -- (-1.275,1.319) -- (-1.237,1.298) -- (-1.200,1.273) -- (-1.162,1.245) -- (-1.125,1.214) -- (-1.089,1.180) -- (-1.053,1.144) -- (-1.019,1.105) -- (-0.986,1.064) -- (-0.955,1.021) -- (-0.925,0.977) -- (-0.897,0.931) -- (-0.872,0.885) -- (-0.848,0.837) -- (-0.827,0.789) -- (-0.809,0.741) -- (-0.793,0.693) -- (-0.780,0.646) -- (-0.770,0.599) -- (-0.763,0.553) -- (-0.758,0.509) -- (-0.757,0.466) -- (-0.758,0.425) -- (-0.763,0.386) -- (-0.770,0.349) -- (-0.780,0.315) -- (-0.793,0.283) -- (-0.809,0.255) -- (-0.827,0.229) -- (-0.848,0.207) -- (-0.872,0.189) -- (-0.897,0.174);
\draw[mesh] (-1.992,1.023) -- (-1.963,1.046) -- (-1.932,1.066) -- (-1.898,1.082) -- (-1.862,1.095) -- (-1.823,1.105) -- (-1.783,1.111) -- (-1.741,1.114) -- (-1.698,1.113) -- (-1.653,1.109) -- (-1.608,1.101) -- (-1.562,1.089) -- (-1.516,1.075) -- (-1.469,1.056) -- (-1.423,1.035) -- (-1.378,1.011) -- (-1.334,0.984) -- (-1.290,0.954) -- (-1.248,0.921) -- (-1.208,0.886) -- (-1.169,0.849) -- (-1.133,0.810) -- (-1.099,0.770) -- (-1.068,0.728) -- (-1.039,0.684) -- (-1.013,0.640) -- (-0.991,0.595) -- (-0.971,0.550) -- (-0.955,0.505) -- (-0.943,0.460) -- (-0.934,0.416) -- (-0.928,0.372) -- (-0.927,0.329) -- (-0.928,0.288) -- (-0.934,0.248) -- (-0.943,0.210) -- (-0.955,0.174) -- (-0.971,0.140) -- (-0.991,0.109) -- (-1.013,0.080) -- (-1.039,0.054);
\draw[mesh] (-2.297,0.584) -- (-2.276,0.619) -- (-2.250,0.651) -- (-2.222,0.681) -- (-2.190,0.709) -- (-2.155,0.734) -- (-2.117,0.757) -- (-2.076,0.776) -- (-2.034,0.792) -- (-1.989,0.805) -- (-1.942,0.815) -- (-1.893,0.822) -- (-1.844,0.825) -- (-1.793,0.825) -- (-1.742,0.822) -- (-1.690,0.815) -- (-1.639,0.805) -- (-1.588,0.791) -- (-1.537,0.775) -- (-1.487,0.755) -- (-1.439,0.732) -- (-1.392,0.707) -- (-1.347,0.679) -- (-1.304,0.649) -- (-1.264,0.616) -- (-1.226,0.581) -- (-1.191,0.544) -- (-1.159,0.506) -- (-1.130,0.466) -- (-1.105,0.426) -- (-1.084,0.384) -- (-1.066,0.341) -- (-1.052,0.299) -- (-1.042,0.256) -- (-1.036,0.213) -- (-1.034,0.170) -- (-1.036,0.129) -- (-1.042,0.088) -- (-1.052,0.048) -- (-1.066,0.010) -- (-1.084,-0.026);
\draw[mesh] (-2.428,0.042) -- (-2.422,0.084) -- (-2.411,0.125) -- (-2.397,0.166) -- (-2.378,0.205) -- (-2.356,0.243) -- (-2.330,0.280) -- (-2.300,0.315) -- (-2.267,0.348) -- (-2.231,0.379) -- (-2.192,0.407) -- (-2.150,0.434) -- (-2.105,0.457) -- (-2.059,0.477) -- (-2.010,0.495) -- (-1.960,0.510) -- (-1.909,0.521) -- (-1.856,0.529) -- (-1.803,0.534) -- (-1.750,0.536) -- (-1.697,0.534) -- (-1.644,0.529) -- (-1.591,0.521) -- (-1.540,0.510) -- (-1.490,0.495) -- (-1.441,0.477) -- (-1.395,0.457) -- (-1.350,0.434) -- (-1.308,0.407) -- (-1.269,0.379) -- (-1.233,0.348) -- (-1.200,0.315) -- (-1.170,0.280) -- (-1.144,0.243) -- (-1.122,0.205) -- (-1.103,0.166) -- (-1.089,0.125) -- (-1.078,0.084) -- (-1.072,0.042) -- (-1.070,0.000) -- (-1.072,-0.042);
\draw[mesh] (-2.347,-0.387) -- (-2.345,-0.345) -- (-2.339,-0.302) -- (-2.329,-0.259) -- (-2.315,-0.216) -- (-2.297,-0.174) -- (-2.276,-0.132) -- (-2.250,-0.091) -- (-2.222,-0.052) -- (-2.190,-0.013) -- (-2.155,0.023) -- (-2.117,0.058) -- (-2.076,0.091) -- (-2.034,0.121) -- (-1.989,0.149) -- (-1.942,0.175) -- (-1.893,0.197) -- (-1.844,0.217) -- (-1.793,0.233) -- (-1.742,0.247) -- (-1.690,0.257) -- (-1.639,0.264) -- (-1.588,0.267) -- (-1.537,0.267) -- (-1.487,0.264) -- (-1.439,0.258) -- (-1.392,0.248) -- (-1.347,0.235) -- (-1.304,0.218) -- (-1.264,0.199) -- (-1.226,0.177) -- (-1.191,0.152) -- (-1.159,0.124) -- (-1.130,0.094) -- (-1.105,0.061) -- (-1.084,0.026) -- (-1.066,-0.010);
\draw[mesh] (-2.315,-0.547) -- (-2.329,-0.509) -- (-2.339,-0.470) -- (-2.345,-0.429) -- (-2.347,-0.387);
\draw[mesh] (-2.104,-0.748) -- (-2.103,-0.705) -- (-2.097,-0.662) -- (-2.088,-0.617) -- (-2.076,-0.572) -- (-2.060,-0.527) -- (-2.040,-0.482) -- (-2.018,-0.437) -- (-1.992,-0.393) -- (-1.963,-0.350) -- (-1.932,-0.308) -- (-1.898,-0.267) -- (-1.862,-0.228) -- (-1.823,-0.191) -- (-1.783,-0.156) -- (-1.741,-0.124) -- (-1.698,-0.094) -- (-1.653,-0.067) -- (-1.608,-0.042) -- (-1.562,-0.021) -- (-1.516,-0.003) -- (-1.469,0.012) -- (-1.423,0.023) -- (-1.378,0.031) -- (-1.334,0.036) -- (-1.290,0.036) -- (-1.248,0.034) -- (-1.208,0.028) -- (-1.169,0.018) -- (-1.133,0.005) -- (-1.099,-0.012) -- (-1.068,-0.032) -- (-1.039,-0.054) -- (-1.013,-0.080);
\draw[mesh] (-2.018,-0.997) -- (-2.040,-0.968) -- (-2.060,-0.937) -- (-2.076,-0.903) -- (-2.088,-0.867) -- (-2.097,-0.829) -- (-2.103,-0.789) -- (-2.104,-0.748);
\draw[mesh] (-1.718,-1.058) -- (-1.717,-1.015) -- (-1.712,-0.970) -- (-1.705,-0.925) -- (-1.695,-0.878) -- (-1.682,-0.830) -- (-1.666,-0.782) -- (-1.647,-0.734) -- (-1.626,-0.686) -- (-1.603,-0.639) -- (-1.577,-0.592) -- (-1.550,-0.547) -- (-1.520,-0.502) -- (-1.489,-0.460) -- (-1.456,-0.419) -- (-1.421,-0.380) -- (-1.386,-0.344) -- (-1.350,-0.310) -- (-1.313,-0.279) -- (-1.275,-0.251) -- (-1.237,-0.226) -- (-1.200,-0.204) -- (-1.162,-0.186) -- (-1.125,-0.172) -- (-1.089,-0.161) -- (-1.053,-0.153) -- (-1.019,-0.150) -- (-0.986,-0.150) -- (-0.955,-0.154) -- (-0.925,-0.162) -- (-0.897,-0.174) -- (-0.872,-0.189);
\draw[mesh] (-1.603,-1.335) -- (-1.626,-1.316) -- (-1.647,-1.294) -- (-1.666,-1.269) -- (-1.682,-1.240) -- (-1.695,-1.209) -- (-1.705,-1.175) -- (-1.712,-1.138) -- (-1.717,-1.099) -- (-1.718,-1.058);
\draw[mesh] (-1.215,-1.296) -- (-1.214,-1.252) -- (-1.211,-1.207) -- (-1.206,-1.161) -- (-1.198,-1.112) -- (-1.189,-1.063) -- (-1.178,-1.013) -- (-1.165,-0.962) -- (-1.150,-0.911) -- (-1.134,-0.861) -- (-1.115,-0.811) -- (-1.096,-0.761) -- (-1.075,-0.713) -- (-1.053,-0.666) -- (-1.029,-0.620) -- (-1.005,-0.577) -- (-0.980,-0.535) -- (-0.954,-0.497) -- (-0.928,-0.461) -- (-0.902,-0.427) -- (-0.875,-0.397) -- (-0.848,-0.370) -- (-0.822,-0.347) -- (-0.796,-0.327) -- (-0.770,-0.311) -- (-0.745,-0.299) -- (-0.721,-0.291) -- (-0.697,-0.287) -- (-0.675,-0.286) -- (-0.654,-0.290) -- (-0.635,-0.298);
\draw[mesh] (-1.115,-1.568) -- (-1.134,-1.557) -- (-1.150,-1.541) -- (-1.165,-1.522) -- (-1.178,-1.499) -- (-1.189,-1.473) -- (-1.198,-1.443) -- (-1.206,-1.411) -- (-1.211,-1.375) -- (-1.214,-1.337) -- (-1.215,-1.296);
\draw[mesh] (-0.629,-1.445) -- (-0.628,-1.402) -- (-0.627,-1.356) -- (-0.624,-1.309) -- (-0.620,-1.260) -- (-0.616,-1.209) -- (-0.610,-1.158) -- (-0.603,-1.106) -- (-0.595,-1.053) -- (-0.587,-1.000) -- (-0.577,-0.948) -- (-0.567,-0.896) -- (-0.556,-0.845) -- (-0.545,-0.795) -- (-0.533,-0.747) -- (-0.520,-0.700) -- (-0.507,-0.656) -- (-0.494,-0.614) -- (-0.480,-0.575) -- (-0.467,-0.538) -- (-0.453,-0.505) -- (-0.439,-0.475) -- (-0.425,-0.448) -- (-0.412,-0.425) -- (-0.399,-0.406) -- (-0.386,-0.391) -- (-0.373,-0.380) -- (-0.361,-0.373) -- (-0.349,-0.369) -- (-0.339,-0.371);
\draw[mesh] (-0.567,-1.711) -- (-0.577,-1.706) -- (-0.587,-1.696) -- (-0.595,-1.683) -- (-0.603,-1.665) -- (-0.610,-1.644) -- (-0.616,-1.619) -- (-0.620,-1.591) -- (-0.624,-1.559) -- (-0.627,-1.524) -- (-0.628,-1.486) -- (-0.629,-1.445);
\draw[mesh] (-0.000,-1.496) -- (-0.000,-1.453) -- (-0.000,-1.407) -- (-0.000,-1.359) -- (-0.000,-1.310) -- (-0.000,-1.259) -- (-0.000,-1.207) -- (-0.000,-1.154) -- (-0.000,-1.101) -- (-0.000,-1.048) -- (-0.000,-0.995) -- (-0.000,-0.942) -- (-0.000,-0.890) -- (-0.000,-0.839) -- (-0.000,-0.790) -- (-0.000,-0.743) -- (-0.000,-0.697) -- (-0.000,-0.654) -- (-0.000,-0.614) -- (-0.000,-0.576) -- (-0.000,-0.542) -- (-0.000,-0.510) -- (-0.000,-0.483) -- (-0.000,-0.459) -- (-0.000,-0.438) -- (-0.000,-0.422) -- (-0.000,-0.410) -- (-0.000,-0.402) -- (-0.000,-0.398) -- (-0.000,-0.398);
\draw[mesh] (-0.000,-1.757) -- (-0.000,-1.752) -- (-0.000,-1.744) -- (-0.000,-1.731) -- (-0.000,-1.714) -- (-0.000,-1.694) -- (-0.000,-1.669) -- (-0.000,-1.641) -- (-0.000,-1.610) -- (-0.000,-1.575) -- (-0.000,-1.537) -- (-0.000,-1.496);
\draw[mesh] (0.629,-1.445) -- (0.628,-1.402) -- (0.627,-1.356) -- (0.624,-1.309) -- (0.620,-1.260) -- (0.616,-1.209) -- (0.610,-1.158) -- (0.603,-1.106) -- (0.595,-1.053) -- (0.587,-1.000) -- (0.577,-0.948) -- (0.567,-0.896) -- (0.556,-0.845) -- (0.545,-0.795) -- (0.533,-0.747) -- (0.520,-0.700) -- (0.507,-0.656) -- (0.494,-0.614) -- (0.480,-0.575) -- (0.467,-0.538) -- (0.453,-0.505) -- (0.439,-0.475) -- (0.425,-0.448) -- (0.412,-0.425) -- (0.399,-0.406) -- (0.386,-0.391) -- (0.373,-0.380) -- (0.361,-0.373) -- (0.349,-0.369) -- (0.339,-0.371);
\draw[mesh] (0.567,-1.711) -- (0.577,-1.706) -- (0.587,-1.696) -- (0.595,-1.683) -- (0.603,-1.665) -- (0.610,-1.644) -- (0.616,-1.619) -- (0.620,-1.591) -- (0.624,-1.559) -- (0.627,-1.524) -- (0.628,-1.486) -- (0.629,-1.445);
\draw[mesh] (1.215,-1.296) -- (1.214,-1.252) -- (1.211,-1.207) -- (1.206,-1.161) -- (1.198,-1.112) -- (1.189,-1.063) -- (1.178,-1.013) -- (1.165,-0.962) -- (1.150,-0.911) -- (1.134,-0.861) -- (1.115,-0.811) -- (1.096,-0.761) -- (1.075,-0.713) -- (1.053,-0.666) -- (1.029,-0.620) -- (1.005,-0.577) -- (0.980,-0.535) -- (0.954,-0.497) -- (0.928,-0.461) -- (0.902,-0.427) -- (0.875,-0.397) -- (0.848,-0.370) -- (0.822,-0.347) -- (0.796,-0.327) -- (0.770,-0.311) -- (0.745,-0.299) -- (0.721,-0.291) -- (0.697,-0.287) -- (0.675,-0.286) -- (0.654,-0.290) -- (0.635,-0.298);
\draw[mesh] (1.115,-1.568) -- (1.134,-1.557) -- (1.150,-1.541) -- (1.165,-1.522) -- (1.178,-1.499) -- (1.189,-1.473) -- (1.198,-1.443) -- (1.206,-1.411) -- (1.211,-1.375) -- (1.214,-1.337) -- (1.215,-1.296);
\draw[mesh] (1.718,-1.058) -- (1.717,-1.015) -- (1.712,-0.970) -- (1.705,-0.925) -- (1.695,-0.878) -- (1.682,-0.830) -- (1.666,-0.782) -- (1.647,-0.734) -- (1.626,-0.686) -- (1.603,-0.639) -- (1.577,-0.592) -- (1.550,-0.547) -- (1.520,-0.502) -- (1.489,-0.460) -- (1.456,-0.419) -- (1.421,-0.380) -- (1.386,-0.344) -- (1.350,-0.310) -- (1.313,-0.279) -- (1.275,-0.251) -- (1.237,-0.226) -- (1.200,-0.204) -- (1.162,-0.186) -- (1.125,-0.172) -- (1.089,-0.161) -- (1.053,-0.153) -- (1.019,-0.150) -- (0.986,-0.150) -- (0.955,-0.154) -- (0.925,-0.162) -- (0.897,-0.174) -- (0.872,-0.189);
\draw[mesh] (1.603,-1.335) -- (1.626,-1.316) -- (1.647,-1.294) -- (1.666,-1.269) -- (1.682,-1.240) -- (1.695,-1.209) -- (1.705,-1.175) -- (1.712,-1.138) -- (1.717,-1.099) -- (1.718,-1.058);
\draw[mesh] (2.104,-0.748) -- (2.103,-0.705) -- (2.097,-0.662) -- (2.088,-0.617) -- (2.076,-0.572) -- (2.060,-0.527) -- (2.040,-0.482) -- (2.018,-0.437) -- (1.992,-0.393) -- (1.963,-0.350) -- (1.932,-0.308) -- (1.898,-0.267) -- (1.862,-0.228) -- (1.823,-0.191) -- (1.783,-0.156) -- (1.741,-0.124) -- (1.698,-0.094) -- (1.653,-0.067) -- (1.608,-0.042) -- (1.562,-0.021) -- (1.516,-0.003) -- (1.469,0.012) -- (1.423,0.023) -- (1.378,0.031) -- (1.334,0.036) -- (1.290,0.036) -- (1.248,0.034) -- (1.208,0.028) -- (1.169,0.018) -- (1.133,0.005) -- (1.099,-0.012) -- (1.068,-0.032) -- (1.039,-0.054) -- (1.013,-0.080);
\draw[mesh] (2.018,-0.997) -- (2.040,-0.968) -- (2.060,-0.937) -- (2.076,-0.903) -- (2.088,-0.867) -- (2.097,-0.829) -- (2.103,-0.789) -- (2.104,-0.748);
\draw[mesh] (2.347,-0.387) -- (2.345,-0.345) -- (2.339,-0.302) -- (2.329,-0.259) -- (2.315,-0.216) -- (2.297,-0.174) -- (2.276,-0.132) -- (2.250,-0.091) -- (2.222,-0.052) -- (2.190,-0.013) -- (2.155,0.023) -- (2.117,0.058) -- (2.076,0.091) -- (2.034,0.121) -- (1.989,0.149) -- (1.942,0.175) -- (1.893,0.197) -- (1.844,0.217) -- (1.793,0.233) -- (1.742,0.247) -- (1.690,0.257) -- (1.639,0.264) -- (1.588,0.267) -- (1.537,0.267) -- (1.487,0.264) -- (1.439,0.258) -- (1.392,0.248) -- (1.347,0.235) -- (1.304,0.218) -- (1.264,0.199) -- (1.226,0.177) -- (1.191,0.152) -- (1.159,0.124) -- (1.130,0.094) -- (1.105,0.061) -- (1.084,0.026) -- (1.066,-0.010);
\draw[mesh] (2.315,-0.547) -- (2.329,-0.509) -- (2.339,-0.470) -- (2.345,-0.429) -- (2.347,-0.387);
\draw[mesh] (-2.428,-0.059) -- (-2.423,-0.117) -- (-2.413,-0.176) -- (-2.400,-0.234) -- (-2.383,-0.292) -- (-2.363,-0.349) -- (-2.339,-0.406) -- (-2.311,-0.462) -- (-2.280,-0.518) -- (-2.245,-0.573) -- (-2.207,-0.626) -- (-2.165,-0.679) -- (-2.120,-0.731) -- (-2.072,-0.782) -- (-2.020,-0.831) -- (-1.966,-0.879) -- (-1.908,-0.926) -- (-1.848,-0.972) -- (-1.784,-1.016) -- (-1.718,-1.058) -- (-1.649,-1.099) -- (-1.578,-1.138) -- (-1.504,-1.175) -- (-1.428,-1.210) -- (-1.350,-1.244) -- (-1.270,-1.276) -- (-1.187,-1.305) -- (-1.103,-1.333) -- (-1.017,-1.359) -- (-0.930,-1.382) -- (-0.841,-1.404) -- (-0.751,-1.423) -- (-0.660,-1.440) -- (-0.567,-1.455) -- (-0.474,-1.467) -- (-0.380,-1.478) -- (-0.286,-1.486) -- (-0.191,-1.491) -- (-0.095,-1.495) -- (-0.000,-1.496) -- (0.095,-1.495) -- (0.191,-1.491) -- (0.286,-1.486) -- (0.380,-1.478) -- (0.474,-1.467) -- (0.567,-1.455) -- (0.660,-1.440) -- (0.751,-1.423) -- (0.841,-1.404) -- (0.930,-1.382) -- (1.017,-1.359) -- (1.103,-1.333) -- (1.187,-1.305) -- (1.270,-1.276) -- (1.350,-1.244) -- (1.428,-1.210) -- (1.504,-1.175) -- (1.578,-1.138) -- (1.649,-1.099) -- (1.718,-1.058) -- (1.784,-1.016) -- (1.848,-0.972) -- (1.908,-0.926) -- (1.966,-0.879) -- (2.020,-0.831) -- (2.072,-0.782) -- (2.120,-0.731) -- (2.165,-0.679) -- (2.207,-0.626) -- (2.245,-0.573) -- (2.280,-0.518) -- (2.311,-0.462) -- (2.339,-0.406) -- (2.363,-0.349) -- (2.383,-0.292) -- (2.400,-0.234) -- (2.413,-0.176) -- (2.423,-0.117) -- (2.428,-0.059) -- (2.430,-0.000);
\draw[mesh] (2.231,0.379) -- (2.229,0.433) -- (2.224,0.487) -- (2.215,0.540) -- (2.203,0.594) -- (2.188,0.647) -- (2.169,0.700) -- (2.147,0.752) -- (2.122,0.803) -- (2.093,0.854) -- (2.061,0.904) -- (2.026,0.954) -- (1.988,1.002) -- (1.946,1.050) -- (1.902,1.097) -- (1.855,1.142) -- (1.805,1.186) -- (1.752,1.229) -- (1.696,1.271) -- (1.638,1.311) -- (1.577,1.350) -- (1.514,1.387) -- (1.449,1.423) -- (1.381,1.457);
\draw[mesh] (-1.449,1.423) -- (-1.514,1.387) -- (-1.577,1.350) -- (-1.638,1.311) -- (-1.696,1.271) -- (-1.752,1.229) -- (-1.805,1.186) -- (-1.855,1.142) -- (-1.902,1.097) -- (-1.946,1.050) -- (-1.988,1.002) -- (-2.026,0.954) -- (-2.061,0.904) -- (-2.093,0.854) -- (-2.122,0.803) -- (-2.147,0.752) -- (-2.169,0.700) -- (-2.188,0.647) -- (-2.203,0.594) -- (-2.215,0.540) -- (-2.224,0.487) -- (-2.229,0.433) -- (-2.231,0.379) -- (-2.229,0.325) -- (-2.224,0.271) -- (-2.215,0.217) -- (-2.203,0.164) -- (-2.188,0.111) -- (-2.169,0.058) -- (-2.147,0.006) -- (-2.122,-0.046) -- (-2.093,-0.096) -- (-2.061,-0.147) -- (-2.026,-0.196) -- (-1.988,-0.245) -- (-1.946,-0.292) -- (-1.902,-0.339) -- (-1.855,-0.384) -- (-1.805,-0.428) -- (-1.752,-0.471) -- (-1.696,-0.513) -- (-1.638,-0.553) -- (-1.577,-0.592) -- (-1.514,-0.630) -- (-1.449,-0.665) -- (-1.381,-0.700) -- (-1.311,-0.732) -- (-1.239,-0.763) -- (-1.166,-0.792) -- (-1.090,-0.819) -- (-1.013,-0.845) -- (-0.934,-0.868) -- (-0.854,-0.890) -- (-0.772,-0.910) -- (-0.689,-0.927) -- (-0.606,-0.943) -- (-0.521,-0.957) -- (-0.435,-0.968) -- (-0.349,-0.978) -- (-0.262,-0.985) -- (-0.175,-0.990) -- (-0.088,-0.993) -- (-0.000,-0.995) -- (0.088,-0.993) -- (0.175,-0.990) -- (0.262,-0.985) -- (0.349,-0.978) -- (0.435,-0.968) -- (0.521,-0.957) -- (0.606,-0.943) -- (0.689,-0.927) -- (0.772,-0.910) -- (0.854,-0.890) -- (0.934,-0.868) -- (1.013,-0.845) -- (1.090,-0.819) -- (1.166,-0.792) -- (1.239,-0.763) -- (1.311,-0.732) -- (1.381,-0.700) -- (1.449,-0.665) -- (1.514,-0.630) -- (1.577,-0.592) -- (1.638,-0.553) -- (1.696,-0.513) -- (1.752,-0.471) -- (1.805,-0.428) -- (1.855,-0.384) -- (1.902,-0.339) -- (1.946,-0.292) -- (1.988,-0.245) -- (2.026,-0.196) -- (2.061,-0.147) -- (2.093,-0.096) -- (2.122,-0.046) -- (2.147,0.006) -- (2.169,0.058) -- (2.188,0.111) -- (2.203,0.164) -- (2.215,0.217) -- (2.224,0.271) -- (2.229,0.325) -- (2.231,0.379);
\draw[mesh] (1.750,0.536) -- (1.749,0.578) -- (1.745,0.620) -- (1.738,0.662) -- (1.728,0.704) -- (1.716,0.746) -- (1.702,0.787) -- (1.684,0.828) -- (1.664,0.869) -- (1.642,0.909) -- (1.617,0.948) -- (1.589,0.987) -- (1.559,1.025) -- (1.527,1.062) -- (1.492,1.099) -- (1.455,1.134) -- (1.416,1.169) -- (1.374,1.203) -- (1.331,1.236) -- (1.285,1.267) -- (1.237,1.298) -- (1.188,1.327) -- (1.137,1.355) -- (1.083,1.382) -- (1.029,1.407) -- (0.972,1.432) -- (0.914,1.454) -- (0.855,1.476) -- (0.794,1.496) -- (0.733,1.514) -- (0.670,1.531) -- (0.606,1.547) -- (0.541,1.561) -- (0.475,1.573) -- (0.409,1.583) -- (0.341,1.593) -- (0.274,1.600) -- (0.206,1.606) -- (0.137,1.610) -- (0.069,1.612) -- (0.000,1.613) -- (-0.069,1.612) -- (-0.137,1.610) -- (-0.206,1.606) -- (-0.274,1.600) -- (-0.341,1.593) -- (-0.409,1.583) -- (-0.475,1.573) -- (-0.541,1.561) -- (-0.606,1.547) -- (-0.670,1.531) -- (-0.733,1.514) -- (-0.794,1.496) -- (-0.855,1.476) -- (-0.914,1.454) -- (-0.972,1.432) -- (-1.029,1.407) -- (-1.083,1.382) -- (-1.137,1.355) -- (-1.188,1.327) -- (-1.237,1.298) -- (-1.285,1.267) -- (-1.331,1.236) -- (-1.374,1.203) -- (-1.416,1.169) -- (-1.455,1.134) -- (-1.492,1.099) -- (-1.527,1.062) -- (-1.559,1.025) -- (-1.589,0.987) -- (-1.617,0.948) -- (-1.642,0.909) -- (-1.664,0.869) -- (-1.684,0.828) -- (-1.702,0.787) -- (-1.716,0.746) -- (-1.728,0.704) -- (-1.738,0.662) -- (-1.745,0.620) -- (-1.749,0.578) -- (-1.750,0.536) -- (-1.749,0.494) -- (-1.745,0.451) -- (-1.738,0.409) -- (-1.728,0.367) -- (-1.716,0.326) -- (-1.702,0.284) -- (-1.684,0.243) -- (-1.664,0.203) -- (-1.642,0.163) -- (-1.617,0.124) -- (-1.589,0.085) -- (-1.559,0.047) -- (-1.527,0.009) -- (-1.492,-0.027) -- (-1.455,-0.063) -- (-1.416,-0.097) -- (-1.374,-0.131) -- (-1.331,-0.164) -- (-1.285,-0.195) -- (-1.237,-0.226) -- (-1.188,-0.255) -- (-1.137,-0.283) -- (-1.083,-0.310) -- (-1.029,-0.336) -- (-0.972,-0.360) -- (-0.914,-0.383) -- (-0.855,-0.404) -- (-0.794,-0.424) -- (-0.733,-0.443) -- (-0.670,-0.460) -- (-0.606,-0.475) -- (-0.541,-0.489) -- (-0.475,-0.501) -- (-0.409,-0.512) -- (-0.341,-0.521) -- (-0.274,-0.528) -- (-0.206,-0.534) -- (-0.137,-0.538) -- (-0.069,-0.541) -- (-0.000,-0.542) -- (0.069,-0.541) -- (0.137,-0.538) -- (0.206,-0.534) -- (0.274,-0.528) -- (0.341,-0.521) -- (0.409,-0.512) -- (0.475,-0.501) -- (0.541,-0.489) -- (0.606,-0.475) -- (0.670,-0.460) -- (0.733,-0.443) -- (0.794,-0.424) -- (0.855,-0.404) -- (0.914,-0.383) -- (0.972,-0.360) -- (1.029,-0.336) -- (1.083,-0.310) -- (1.137,-0.283) -- (1.188,-0.255) -- (1.237,-0.226) -- (1.285,-0.195) -- (1.331,-0.164) -- (1.374,-0.131) -- (1.416,-0.097) -- (1.455,-0.063) -- (1.492,-0.027) -- (1.527,0.009) -- (1.559,0.047) -- (1.589,0.085) -- (1.617,0.124) -- (1.642,0.163) -- (1.664,0.203) -- (1.684,0.243) -- (1.702,0.284) -- (1.716,0.326) -- (1.728,0.367) -- (1.738,0.409) -- (1.745,0.451) -- (1.749,0.494) -- (1.750,0.536);
\draw[mesh] (1.269,0.379) -- (1.268,0.410) -- (1.265,0.440) -- (1.260,0.471) -- (1.254,0.501) -- (1.245,0.531) -- (1.234,0.561) -- (1.222,0.591) -- (1.207,0.620) -- (1.191,0.649) -- (1.173,0.678) -- (1.153,0.706) -- (1.131,0.734) -- (1.107,0.761) -- (1.082,0.787) -- (1.055,0.813) -- (1.027,0.838) -- (0.997,0.863) -- (0.965,0.886) -- (0.932,0.909) -- (0.897,0.931) -- (0.862,0.953) -- (0.824,0.973) -- (0.786,0.993) -- (0.746,1.011) -- (0.705,1.029) -- (0.663,1.045) -- (0.620,1.061) -- (0.576,1.075) -- (0.531,1.089) -- (0.486,1.101) -- (0.439,1.112) -- (0.392,1.122) -- (0.345,1.131) -- (0.296,1.139) -- (0.248,1.145) -- (0.199,1.151) -- (0.149,1.155) -- (0.100,1.158) -- (0.050,1.160) -- (0.000,1.160) -- (-0.050,1.160) -- (-0.100,1.158) -- (-0.149,1.155) -- (-0.199,1.151) -- (-0.248,1.145) -- (-0.296,1.139) -- (-0.345,1.131) -- (-0.392,1.122) -- (-0.439,1.112) -- (-0.486,1.101) -- (-0.531,1.089) -- (-0.576,1.075) -- (-0.620,1.061) -- (-0.663,1.045) -- (-0.705,1.029) -- (-0.746,1.011) -- (-0.786,0.993) -- (-0.824,0.973) -- (-0.862,0.953) -- (-0.897,0.931) -- (-0.932,0.909) -- (-0.965,0.886) -- (-0.997,0.863) -- (-1.027,0.838) -- (-1.055,0.813) -- (-1.082,0.787) -- (-1.107,0.761) -- (-1.131,0.734) -- (-1.153,0.706) -- (-1.173,0.678) -- (-1.191,0.649) -- (-1.207,0.620) -- (-1.222,0.591) -- (-1.234,0.561) -- (-1.245,0.531) -- (-1.254,0.501) -- (-1.260,0.471) -- (-1.265,0.440) -- (-1.268,0.410) -- (-1.269,0.379) -- (-1.268,0.348) -- (-1.265,0.318) -- (-1.260,0.287) -- (-1.254,0.257) -- (-1.245,0.226) -- (-1.234,0.196) -- (-1.222,0.167) -- (-1.207,0.137) -- (-1.191,0.108) -- (-1.173,0.080) -- (-1.153,0.052) -- (-1.131,0.024) -- (-1.107,-0.003) -- (-1.082,-0.029) -- (-1.055,-0.055) -- (-1.027,-0.080) -- (-0.997,-0.105) -- (-0.965,-0.129) -- (-0.932,-0.151) -- (-0.897,-0.174) -- (-0.862,-0.195) -- (-0.824,-0.215) -- (-0.786,-0.235);
\draw[mesh] (0.824,-0.215) -- (0.862,-0.195) -- (0.897,-0.174) -- (0.932,-0.151) -- (0.965,-0.129) -- (0.997,-0.105) -- (1.027,-0.080) -- (1.055,-0.055) -- (1.082,-0.029) -- (1.107,-0.003) -- (1.131,0.024) -- (1.153,0.052) -- (1.173,0.080) -- (1.191,0.108) -- (1.207,0.137) -- (1.222,0.167) -- (1.234,0.196) -- (1.245,0.226) -- (1.254,0.257) -- (1.260,0.287) -- (1.265,0.318) -- (1.268,0.348) -- (1.269,0.379);
\draw[mesh] (1.070,0.000) -- (1.069,0.026) -- (1.067,0.052) -- (1.063,0.077) -- (1.057,0.103) -- (1.049,0.129) -- (1.040,0.154) -- (1.030,0.179) -- (1.018,0.204) -- (1.004,0.228) -- (0.989,0.252) -- (0.972,0.276) -- (0.953,0.299) -- (0.934,0.322) -- (0.912,0.344) -- (0.890,0.366) -- (0.866,0.387) -- (0.840,0.408) -- (0.814,0.428) -- (0.786,0.447) -- (0.757,0.466) -- (0.726,0.484) -- (0.695,0.501) -- (0.662,0.517) -- (0.629,0.533) -- (0.594,0.548) -- (0.559,0.562) -- (0.523,0.575) -- (0.486,0.587) -- (0.448,0.598) -- (0.409,0.609) -- (0.370,0.618) -- (0.331,0.627) -- (0.290,0.634) -- (0.250,0.641) -- (0.209,0.646) -- (0.167,0.651) -- (0.126,0.654) -- (0.084,0.657) -- (0.042,0.658) -- (0.000,0.659) -- (-0.042,0.658) -- (-0.084,0.657) -- (-0.126,0.654) -- (-0.167,0.651) -- (-0.209,0.646) -- (-0.250,0.641) -- (-0.290,0.634) -- (-0.331,0.627) -- (-0.370,0.618) -- (-0.409,0.609) -- (-0.448,0.598) -- (-0.486,0.587) -- (-0.523,0.575) -- (-0.559,0.562) -- (-0.594,0.548) -- (-0.629,0.533) -- (-0.662,0.517) -- (-0.695,0.501) -- (-0.726,0.484) -- (-0.757,0.466) -- (-0.786,0.447) -- (-0.814,0.428) -- (-0.840,0.408) -- (-0.866,0.387) -- (-0.890,0.366) -- (-0.912,0.344) -- (-0.934,0.322) -- (-0.953,0.299) -- (-0.972,0.276) -- (-0.989,0.252) -- (-1.004,0.228) -- (-1.018,0.204) -- (-1.030,0.179) -- (-1.040,0.154) -- (-1.049,0.129) -- (-1.057,0.103) -- (-1.063,0.077) -- (-1.067,0.052) -- (-1.069,0.026) -- (-1.070,0.000) -- (-1.069,-0.026);
\draw[mesh] (0.786,0.235) -- (0.746,0.253) -- (0.705,0.271) -- (0.663,0.287) -- (0.620,0.303) -- (0.576,0.317) -- (0.531,0.331) -- (0.486,0.343) -- (0.439,0.354) -- (0.392,0.364) -- (0.345,0.373) -- (0.296,0.381) -- (0.248,0.387) -- (0.199,0.393) -- (0.149,0.397) -- (0.100,0.400) -- (0.050,0.402) -- (0.000,0.402) -- (-0.050,0.402) -- (-0.100,0.400) -- (-0.149,0.397) -- (-0.199,0.393) -- (-0.248,0.387) -- (-0.296,0.381) -- (-0.345,0.373) -- (-0.392,0.364) -- (-0.439,0.354) -- (-0.486,0.343) -- (-0.531,0.331) -- (-0.576,0.317) -- (-0.620,0.303) -- (-0.663,0.287) -- (-0.705,0.271) -- (-0.746,0.253) -- (-0.786,0.235) -- (-0.824,0.215);
\draw[mesh] (-1.381,-1.457) -- (-1.311,-1.490) -- (-1.239,-1.521) -- (-1.166,-1.550) -- (-1.090,-1.577) -- (-1.013,-1.603) -- (-0.934,-1.626) -- (-0.854,-1.648) -- (-0.772,-1.667) -- (-0.689,-1.685) -- (-0.606,-1.701) -- (-0.521,-1.714) -- (-0.435,-1.726) -- (-0.349,-1.735) -- (-0.262,-1.743) -- (-0.175,-1.748) -- (-0.088,-1.751) -- (-0.000,-1.752) -- (0.088,-1.751) -- (0.175,-1.748) -- (0.262,-1.743) -- (0.349,-1.735) -- (0.435,-1.726) -- (0.521,-1.714) -- (0.606,-1.701) -- (0.689,-1.685) -- (0.772,-1.667) -- (0.854,-1.648) -- (0.934,-1.626) -- (1.013,-1.603) -- (1.090,-1.577) -- (1.166,-1.550) -- (1.239,-1.521) -- (1.311,-1.490) -- (1.381,-1.457) -- (1.449,-1.423);

    % a scan winding at f2/f1 = golden ratio, generated by gen-loops.py
\draw[scanback] (2.430,0.000) -- (2.430,0.019);
\draw[scan] (2.430,0.019) -- (2.429,0.037) -- (2.429,0.056) -- (2.428,0.074) -- (2.427,0.093) -- (2.425,0.111) -- (2.424,0.130) -- (2.422,0.148) -- (2.420,0.167) -- (2.417,0.185) -- (2.414,0.203) -- (2.411,0.222) -- (2.408,0.240) -- (2.405,0.258) -- (2.401,0.276) -- (2.397,0.295) -- (2.393,0.313) -- (2.388,0.331) -- (2.383,0.349) -- (2.379,0.367) -- (2.373,0.385) -- (2.368,0.402) -- (2.362,0.420) -- (2.356,0.438) -- (2.350,0.455) -- (2.343,0.473) -- (2.337,0.490) -- (2.330,0.507) -- (2.323,0.525) -- (2.315,0.542) -- (2.308,0.559) -- (2.300,0.576) -- (2.292,0.592) -- (2.284,0.609) -- (2.275,0.626) -- (2.266,0.642) -- (2.257,0.658) -- (2.248,0.675) -- (2.239,0.691) -- (2.229,0.707) -- (2.220,0.723) -- (2.210,0.738) -- (2.199,0.754) -- (2.189,0.769) -- (2.178,0.785) -- (2.168,0.800) -- (2.157,0.815) -- (2.145,0.830) -- (2.134,0.845) -- (2.122,0.859) -- (2.111,0.874) -- (2.099,0.888) -- (2.087,0.902) -- (2.074,0.916) -- (2.062,0.930) -- (2.049,0.944) -- (2.037,0.957) -- (2.024,0.971) -- (2.011,0.984) -- (1.997,0.997) -- (1.984,1.010) -- (1.970,1.022) -- (1.957,1.035) -- (1.943,1.047) -- (1.929,1.059) -- (1.915,1.071) -- (1.901,1.083) -- (1.886,1.094) -- (1.872,1.106) -- (1.857,1.117) -- (1.843,1.128) -- (1.828,1.139) -- (1.813,1.149) -- (1.798,1.160) -- (1.783,1.170) -- (1.767,1.180) -- (1.752,1.190) -- (1.736,1.200) -- (1.721,1.209) -- (1.705,1.218) -- (1.690,1.228) -- (1.674,1.236) -- (1.658,1.245) -- (1.642,1.254) -- (1.626,1.262) -- (1.610,1.270) -- (1.594,1.278) -- (1.577,1.285) -- (1.561,1.293) -- (1.545,1.300) -- (1.529,1.307) -- (1.512,1.314) -- (1.496,1.321) -- (1.479,1.327) -- (1.463,1.333) -- (1.446,1.339) -- (1.429,1.345) -- (1.413,1.351) -- (1.396,1.356) -- (1.380,1.361) -- (1.363,1.366) -- (1.346,1.371) -- (1.329,1.376) -- (1.313,1.380) -- (1.296,1.384) -- (1.279,1.388) -- (1.263,1.392) -- (1.246,1.395) -- (1.229,1.399) -- (1.213,1.402) -- (1.196,1.405) -- (1.179,1.408) -- (1.163,1.410) -- (1.146,1.412) -- (1.129,1.415) -- (1.113,1.417) -- (1.096,1.418) -- (1.080,1.420) -- (1.063,1.421) -- (1.047,1.422) -- (1.031,1.423) -- (1.014,1.424) -- (0.998,1.425) -- (0.982,1.425) -- (0.966,1.425) -- (0.949,1.425) -- (0.933,1.425) -- (0.917,1.425) -- (0.901,1.424) -- (0.886,1.424) -- (0.870,1.423) -- (0.854,1.422) -- (0.838,1.421) -- (0.823,1.419) -- (0.807,1.418) -- (0.792,1.416) -- (0.776,1.414) -- (0.761,1.412) -- (0.746,1.410) -- (0.731,1.407) -- (0.716,1.404) -- (0.701,1.402) -- (0.686,1.399) -- (0.671,1.396) -- (0.656,1.393) -- (0.642,1.389) -- (0.627,1.386) -- (0.613,1.382) -- (0.598,1.378) -- (0.584,1.374) -- (0.570,1.370) -- (0.556,1.366) -- (0.542,1.361) -- (0.528,1.357) -- (0.514,1.352) -- (0.501,1.347) -- (0.487,1.342) -- (0.474,1.337) -- (0.460,1.332) -- (0.447,1.327) -- (0.434,1.321) -- (0.421,1.316) -- (0.408,1.310) -- (0.395,1.304) -- (0.382,1.298) -- (0.370,1.292) -- (0.357,1.286) -- (0.345,1.280) -- (0.332,1.273) -- (0.320,1.267) -- (0.308,1.260) -- (0.296,1.253) -- (0.284,1.247) -- (0.272,1.240) -- (0.261,1.233) -- (0.249,1.226) -- (0.238,1.219) -- (0.226,1.211) -- (0.215,1.204) -- (0.204,1.197) -- (0.193,1.189) -- (0.182,1.182) -- (0.171,1.174) -- (0.160,1.166) -- (0.149,1.158) -- (0.139,1.151) -- (0.128,1.143) -- (0.118,1.135) -- (0.108,1.127) -- (0.097,1.118) -- (0.087,1.110) -- (0.077,1.102) -- (0.067,1.094) -- (0.057,1.085) -- (0.048,1.077) -- (0.038,1.069) -- (0.028,1.060) -- (0.019,1.052) -- (0.009,1.043) -- (0.000,1.034) -- (-0.009,1.026) -- (-0.019,1.017) -- (-0.028,1.008) -- (-0.037,1.000) -- (-0.046,0.991) -- (-0.055,0.982) -- (-0.064,0.973) -- (-0.072,0.964) -- (-0.081,0.956) -- (-0.090,0.947) -- (-0.098,0.938) -- (-0.107,0.929) -- (-0.116,0.920) -- (-0.124,0.911) -- (-0.132,0.902) -- (-0.141,0.893) -- (-0.149,0.884) -- (-0.157,0.875) -- (-0.166,0.866) -- (-0.174,0.857) -- (-0.182,0.848) -- (-0.190,0.839) -- (-0.198,0.830) -- (-0.206,0.821) -- (-0.214,0.812) -- (-0.222,0.803) -- (-0.230,0.794) -- (-0.238,0.785) -- (-0.246,0.776) -- (-0.254,0.767) -- (-0.261,0.758) -- (-0.269,0.749) -- (-0.277,0.740) -- (-0.285,0.731) -- (-0.293,0.722) -- (-0.300,0.713) -- (-0.308,0.704) -- (-0.316,0.695) -- (-0.324,0.686) -- (-0.332,0.677) -- (-0.339,0.668) -- (-0.347,0.660) -- (-0.355,0.651) -- (-0.363,0.642) -- (-0.370,0.633) -- (-0.378,0.625) -- (-0.386,0.616) -- (-0.394,0.607) -- (-0.402,0.599) -- (-0.410,0.590) -- (-0.418,0.581) -- (-0.425,0.573) -- (-0.433,0.564) -- (-0.441,0.556) -- (-0.449,0.547) -- (-0.457,0.539) -- (-0.465,0.530) -- (-0.474,0.522) -- (-0.482,0.514) -- (-0.490,0.505) -- (-0.498,0.497) -- (-0.506,0.489) -- (-0.515,0.480) -- (-0.523,0.472) -- (-0.531,0.464) -- (-0.540,0.456) -- (-0.548,0.448) -- (-0.557,0.440) -- (-0.565,0.432) -- (-0.574,0.424) -- (-0.583,0.416) -- (-0.591,0.408) -- (-0.600,0.400) -- (-0.609,0.392) -- (-0.618,0.384) -- (-0.627,0.376) -- (-0.636,0.368) -- (-0.645,0.361) -- (-0.654,0.353) -- (-0.663,0.345) -- (-0.672,0.337) -- (-0.682,0.330) -- (-0.691,0.322) -- (-0.701,0.315) -- (-0.710,0.307) -- (-0.720,0.299) -- (-0.729,0.292) -- (-0.739,0.284) -- (-0.749,0.277) -- (-0.759,0.270) -- (-0.769,0.262) -- (-0.778,0.255) -- (-0.789,0.248) -- (-0.799,0.240) -- (-0.809,0.233) -- (-0.819,0.226) -- (-0.829,0.218) -- (-0.840,0.211) -- (-0.850,0.204) -- (-0.861,0.197) -- (-0.871,0.190) -- (-0.882,0.183) -- (-0.893,0.175) -- (-0.904,0.168);
\draw[scanback] (-0.904,0.168) -- (-0.914,0.161) -- (-0.925,0.154) -- (-0.936,0.147) -- (-0.947,0.140) -- (-0.959,0.133) -- (-0.970,0.126) -- (-0.981,0.119) -- (-0.992,0.112) -- (-1.004,0.105) -- (-1.015,0.099) -- (-1.026,0.092) -- (-1.038,0.085) -- (-1.050,0.078) -- (-1.061,0.071) -- (-1.073,0.064) -- (-1.085,0.057) -- (-1.096,0.050) -- (-1.108,0.044) -- (-1.120,0.037) -- (-1.132,0.030) -- (-1.144,0.023) -- (-1.156,0.016) -- (-1.168,0.010) -- (-1.180,0.003) -- (-1.192,-0.004) -- (-1.205,-0.011) -- (-1.217,-0.018) -- (-1.229,-0.024) -- (-1.241,-0.031) -- (-1.254,-0.038) -- (-1.266,-0.045) -- (-1.278,-0.052) -- (-1.291,-0.058) -- (-1.303,-0.065) -- (-1.316,-0.072) -- (-1.328,-0.079) -- (-1.340,-0.085) -- (-1.353,-0.092) -- (-1.365,-0.099) -- (-1.378,-0.106) -- (-1.390,-0.113) -- (-1.403,-0.119) -- (-1.415,-0.126) -- (-1.428,-0.133) -- (-1.440,-0.140) -- (-1.452,-0.147) -- (-1.465,-0.154) -- (-1.477,-0.160) -- (-1.490,-0.167) -- (-1.502,-0.174) -- (-1.514,-0.181) -- (-1.527,-0.188) -- (-1.539,-0.195) -- (-1.551,-0.202) -- (-1.564,-0.209) -- (-1.576,-0.216) -- (-1.588,-0.223) -- (-1.600,-0.229) -- (-1.612,-0.236) -- (-1.624,-0.243) -- (-1.636,-0.250) -- (-1.648,-0.257) -- (-1.660,-0.264) -- (-1.671,-0.271) -- (-1.683,-0.278) -- (-1.695,-0.285) -- (-1.706,-0.292) -- (-1.718,-0.299) -- (-1.729,-0.306) -- (-1.741,-0.313) -- (-1.752,-0.321) -- (-1.763,-0.328) -- (-1.774,-0.335) -- (-1.785,-0.342) -- (-1.796,-0.349) -- (-1.807,-0.356) -- (-1.817,-0.363) -- (-1.828,-0.370) -- (-1.838,-0.377) -- (-1.849,-0.385) -- (-1.859,-0.392) -- (-1.869,-0.399) -- (-1.879,-0.406) -- (-1.889,-0.413) -- (-1.899,-0.420) -- (-1.908,-0.428) -- (-1.918,-0.435) -- (-1.927,-0.442) -- (-1.936,-0.449) -- (-1.945,-0.456) -- (-1.954,-0.463) -- (-1.963,-0.471) -- (-1.972,-0.478) -- (-1.980,-0.485) -- (-1.988,-0.492) -- (-1.996,-0.499) -- (-2.004,-0.507) -- (-2.012,-0.514) -- (-2.020,-0.521) -- (-2.027,-0.528) -- (-2.035,-0.535) -- (-2.042,-0.542) -- (-2.049,-0.550) -- (-2.055,-0.557) -- (-2.062,-0.564) -- (-2.068,-0.571) -- (-2.075,-0.578) -- (-2.081,-0.585) -- (-2.086,-0.592) -- (-2.092,-0.599) -- (-2.097,-0.606) -- (-2.103,-0.613) -- (-2.108,-0.620) -- (-2.112,-0.627) -- (-2.117,-0.634) -- (-2.121,-0.641) -- (-2.126,-0.648) -- (-2.130,-0.655) -- (-2.133,-0.662) -- (-2.137,-0.669) -- (-2.140,-0.676) -- (-2.143,-0.683) -- (-2.146,-0.689) -- (-2.149,-0.696) -- (-2.151,-0.703) -- (-2.153,-0.710) -- (-2.155,-0.716) -- (-2.157,-0.723) -- (-2.159,-0.729) -- (-2.160,-0.736) -- (-2.161,-0.742) -- (-2.162,-0.749) -- (-2.162,-0.755) -- (-2.163,-0.762) -- (-2.163,-0.768) -- (-2.162,-0.774) -- (-2.162,-0.781) -- (-2.161,-0.787) -- (-2.161,-0.793) -- (-2.159,-0.799) -- (-2.158,-0.805) -- (-2.156,-0.811) -- (-2.155,-0.817) -- (-2.152,-0.823) -- (-2.150,-0.829) -- (-2.148,-0.834) -- (-2.145,-0.840) -- (-2.142,-0.846) -- (-2.138,-0.851) -- (-2.135,-0.857) -- (-2.131,-0.862) -- (-2.127,-0.868);
\draw[scan] (-2.127,-0.868) -- (-2.122,-0.873) -- (-2.118,-0.878) -- (-2.113,-0.883) -- (-2.108,-0.888) -- (-2.103,-0.893) -- (-2.097,-0.898) -- (-2.091,-0.903) -- (-2.085,-0.908) -- (-2.079,-0.913) -- (-2.073,-0.917) -- (-2.066,-0.922) -- (-2.059,-0.926) -- (-2.052,-0.931) -- (-2.044,-0.935) -- (-2.036,-0.939) -- (-2.029,-0.943) -- (-2.020,-0.947) -- (-2.012,-0.951) -- (-2.003,-0.955) -- (-1.994,-0.959) -- (-1.985,-0.962) -- (-1.976,-0.966) -- (-1.966,-0.970) -- (-1.957,-0.973) -- (-1.947,-0.976) -- (-1.936,-0.979) -- (-1.926,-0.982) -- (-1.915,-0.985) -- (-1.904,-0.988) -- (-1.893,-0.991) -- (-1.882,-0.994) -- (-1.871,-0.996) -- (-1.859,-0.999) -- (-1.847,-1.001) -- (-1.835,-1.004) -- (-1.823,-1.006) -- (-1.810,-1.008) -- (-1.797,-1.010) -- (-1.784,-1.012) -- (-1.771,-1.013) -- (-1.758,-1.015) -- (-1.745,-1.017) -- (-1.731,-1.018) -- (-1.717,-1.019) -- (-1.703,-1.020) -- (-1.689,-1.021) -- (-1.675,-1.022) -- (-1.660,-1.023) -- (-1.645,-1.024) -- (-1.630,-1.025) -- (-1.615,-1.025) -- (-1.600,-1.026) -- (-1.585,-1.026) -- (-1.570,-1.026) -- (-1.554,-1.026) -- (-1.538,-1.026) -- (-1.522,-1.026) -- (-1.506,-1.026) -- (-1.490,-1.025) -- (-1.474,-1.025) -- (-1.457,-1.024) -- (-1.441,-1.023) -- (-1.424,-1.022) -- (-1.407,-1.022) -- (-1.391,-1.020) -- (-1.374,-1.019) -- (-1.356,-1.018) -- (-1.339,-1.017) -- (-1.322,-1.015) -- (-1.304,-1.013) -- (-1.287,-1.012) -- (-1.269,-1.010) -- (-1.252,-1.008) -- (-1.234,-1.006) -- (-1.216,-1.003) -- (-1.198,-1.001) -- (-1.180,-0.999) -- (-1.162,-0.996) -- (-1.144,-0.994) -- (-1.126,-0.991) -- (-1.108,-0.988) -- (-1.090,-0.985) -- (-1.071,-0.982) -- (-1.053,-0.979) -- (-1.034,-0.976) -- (-1.016,-0.972) -- (-0.998,-0.969) -- (-0.979,-0.965) -- (-0.961,-0.962) -- (-0.942,-0.958) -- (-0.923,-0.954) -- (-0.905,-0.950) -- (-0.886,-0.946) -- (-0.868,-0.942) -- (-0.849,-0.938) -- (-0.830,-0.933) -- (-0.812,-0.929) -- (-0.793,-0.925) -- (-0.775,-0.920) -- (-0.756,-0.915) -- (-0.738,-0.911) -- (-0.719,-0.906) -- (-0.701,-0.901) -- (-0.682,-0.896) -- (-0.664,-0.891) -- (-0.645,-0.886) -- (-0.627,-0.880) -- (-0.609,-0.875) -- (-0.590,-0.870) -- (-0.572,-0.864) -- (-0.554,-0.859) -- (-0.536,-0.853) -- (-0.518,-0.848) -- (-0.500,-0.842) -- (-0.482,-0.836) -- (-0.464,-0.830) -- (-0.446,-0.824) -- (-0.428,-0.819) -- (-0.410,-0.813) -- (-0.393,-0.806) -- (-0.375,-0.800) -- (-0.358,-0.794) -- (-0.340,-0.788) -- (-0.323,-0.782) -- (-0.306,-0.776) -- (-0.289,-0.769) -- (-0.272,-0.763) -- (-0.255,-0.756) -- (-0.238,-0.750) -- (-0.222,-0.744) -- (-0.205,-0.737) -- (-0.189,-0.731) -- (-0.172,-0.724) -- (-0.156,-0.717) -- (-0.140,-0.711) -- (-0.124,-0.704) -- (-0.108,-0.698) -- (-0.092,-0.691) -- (-0.076,-0.684) -- (-0.061,-0.678) -- (-0.045,-0.671) -- (-0.030,-0.664) -- (-0.015,-0.657) -- (-0.000,-0.651) -- (0.015,-0.644) -- (0.030,-0.637) -- (0.044,-0.630) -- (0.059,-0.624) -- (0.073,-0.617) -- (0.087,-0.610) -- (0.101,-0.604) -- (0.115,-0.597) -- (0.129,-0.590) -- (0.143,-0.584) -- (0.156,-0.577) -- (0.170,-0.570) -- (0.183,-0.564) -- (0.196,-0.557) -- (0.209,-0.551) -- (0.222,-0.544) -- (0.234,-0.537) -- (0.247,-0.531) -- (0.259,-0.525) -- (0.271,-0.518) -- (0.283,-0.512) -- (0.295,-0.505) -- (0.307,-0.499) -- (0.318,-0.493) -- (0.330,-0.487) -- (0.341,-0.480) -- (0.352,-0.474) -- (0.363,-0.468) -- (0.374,-0.462) -- (0.384,-0.456) -- (0.395,-0.450) -- (0.405,-0.445) -- (0.416,-0.439) -- (0.426,-0.433) -- (0.436,-0.427) -- (0.445,-0.422) -- (0.455,-0.416) -- (0.464,-0.411) -- (0.474,-0.405) -- (0.483,-0.400) -- (0.492,-0.395) -- (0.501,-0.389) -- (0.510,-0.384) -- (0.518,-0.379) -- (0.527,-0.374) -- (0.535,-0.369) -- (0.543,-0.364) -- (0.551,-0.360) -- (0.559,-0.355) -- (0.567,-0.350) -- (0.575,-0.346) -- (0.582,-0.341) -- (0.589,-0.337) -- (0.597,-0.333) -- (0.604,-0.329) -- (0.611,-0.324) -- (0.618,-0.320) -- (0.624,-0.316) -- (0.631,-0.313) -- (0.638,-0.309) -- (0.644,-0.305) -- (0.650,-0.302) -- (0.656,-0.298) -- (0.662,-0.295) -- (0.668,-0.292) -- (0.674,-0.288) -- (0.680,-0.285) -- (0.686,-0.282) -- (0.691,-0.279) -- (0.696,-0.276) -- (0.702,-0.274) -- (0.707,-0.271) -- (0.712,-0.269) -- (0.717,-0.266) -- (0.722,-0.264) -- (0.727,-0.262) -- (0.732,-0.259) -- (0.736,-0.257) -- (0.741,-0.255);
\draw[scanback] (0.741,-0.255) -- (0.745,-0.254) -- (0.750,-0.252) -- (0.754,-0.250) -- (0.758,-0.249) -- (0.763,-0.247) -- (0.767,-0.246) -- (0.771,-0.244) -- (0.775,-0.243) -- (0.779,-0.242) -- (0.783,-0.241) -- (0.787,-0.240) -- (0.790,-0.239) -- (0.794,-0.239) -- (0.798,-0.238) -- (0.801,-0.237) -- (0.805,-0.237) -- (0.809,-0.236) -- (0.812,-0.236) -- (0.816,-0.236) -- (0.819,-0.236) -- (0.822,-0.236) -- (0.826,-0.236) -- (0.829,-0.236) -- (0.832,-0.236) -- (0.836,-0.236) -- (0.839,-0.237) -- (0.842,-0.237) -- (0.845,-0.238) -- (0.849,-0.238) -- (0.852,-0.239) -- (0.855,-0.240) -- (0.858,-0.241) -- (0.861,-0.242) -- (0.864,-0.242) -- (0.867,-0.243) -- (0.871,-0.245) -- (0.874,-0.246) -- (0.877,-0.247) -- (0.880,-0.248) -- (0.883,-0.250) -- (0.886,-0.251) -- (0.889,-0.252) -- (0.893,-0.254) -- (0.896,-0.255) -- (0.899,-0.257) -- (0.902,-0.259) -- (0.906,-0.260) -- (0.909,-0.262) -- (0.912,-0.264) -- (0.915,-0.266) -- (0.919,-0.268) -- (0.922,-0.270) -- (0.925,-0.271) -- (0.929,-0.273) -- (0.932,-0.275) -- (0.936,-0.277) -- (0.939,-0.279) -- (0.943,-0.282) -- (0.946,-0.284) -- (0.950,-0.286) -- (0.954,-0.288) -- (0.957,-0.290) -- (0.961,-0.292) -- (0.965,-0.294) -- (0.968,-0.297) -- (0.972,-0.299) -- (0.976,-0.301) -- (0.980,-0.303) -- (0.984,-0.305) -- (0.988,-0.308) -- (0.992,-0.310) -- (0.996,-0.312) -- (1.000,-0.314) -- (1.005,-0.316) -- (1.009,-0.318) -- (1.013,-0.321) -- (1.017,-0.323) -- (1.022,-0.325) -- (1.026,-0.327) -- (1.031,-0.329) -- (1.035,-0.331) -- (1.040,-0.333) -- (1.044,-0.335) -- (1.049,-0.337) -- (1.054,-0.339) -- (1.058,-0.341) -- (1.063,-0.342) -- (1.068,-0.344) -- (1.073,-0.346) -- (1.078,-0.347) -- (1.083,-0.349) -- (1.088,-0.351) -- (1.093,-0.352) -- (1.098,-0.353) -- (1.103,-0.355) -- (1.108,-0.356) -- (1.114,-0.357) -- (1.119,-0.359) -- (1.124,-0.360) -- (1.130,-0.361) -- (1.135,-0.362) -- (1.140,-0.363) -- (1.146,-0.363) -- (1.151,-0.364) -- (1.157,-0.365) -- (1.163,-0.365) -- (1.168,-0.366) -- (1.174,-0.366) -- (1.179,-0.366) -- (1.185,-0.367) -- (1.191,-0.367) -- (1.196,-0.367) -- (1.202,-0.367) -- (1.208,-0.366) -- (1.214,-0.366) -- (1.219,-0.366) -- (1.225,-0.365) -- (1.231,-0.364) -- (1.237,-0.364) -- (1.243,-0.363) -- (1.249,-0.362) -- (1.254,-0.361) -- (1.260,-0.360) -- (1.266,-0.358) -- (1.272,-0.357) -- (1.278,-0.355) -- (1.284,-0.354) -- (1.289,-0.352) -- (1.295,-0.350) -- (1.301,-0.348) -- (1.307,-0.346) -- (1.313,-0.343) -- (1.318,-0.341) -- (1.324,-0.338) -- (1.330,-0.336) -- (1.335,-0.333) -- (1.341,-0.330) -- (1.347,-0.327) -- (1.352,-0.324) -- (1.358,-0.320) -- (1.363,-0.317) -- (1.369,-0.313) -- (1.374,-0.309) -- (1.379,-0.305) -- (1.385,-0.301) -- (1.390,-0.297) -- (1.395,-0.293) -- (1.400,-0.288) -- (1.405,-0.284) -- (1.410,-0.279) -- (1.415,-0.274) -- (1.420,-0.269) -- (1.425,-0.263) -- (1.430,-0.258) -- (1.434,-0.253) -- (1.439,-0.247) -- (1.443,-0.241) -- (1.448,-0.235) -- (1.452,-0.229) -- (1.456,-0.223) -- (1.460,-0.216) -- (1.464,-0.210) -- (1.468,-0.203) -- (1.472,-0.196) -- (1.476,-0.189) -- (1.479,-0.182) -- (1.483,-0.175) -- (1.486,-0.167) -- (1.490,-0.160) -- (1.493,-0.152) -- (1.496,-0.144) -- (1.499,-0.136) -- (1.502,-0.128) -- (1.504,-0.120) -- (1.507,-0.112) -- (1.509,-0.103) -- (1.511,-0.094) -- (1.513,-0.085) -- (1.515,-0.077) -- (1.517,-0.067) -- (1.519,-0.058) -- (1.520,-0.049) -- (1.522,-0.039) -- (1.523,-0.030) -- (1.524,-0.020) -- (1.525,-0.010) -- (1.526,-0.000) -- (1.526,0.010) -- (1.527,0.020) -- (1.527,0.031) -- (1.527,0.041) -- (1.527,0.052) -- (1.527,0.063) -- (1.526,0.074) -- (1.526,0.084) -- (1.525,0.096) -- (1.524,0.107) -- (1.523,0.118) -- (1.522,0.130) -- (1.520,0.141) -- (1.518,0.153) -- (1.516,0.164) -- (1.514,0.176) -- (1.512,0.188) -- (1.510,0.200) -- (1.507,0.212) -- (1.504,0.225) -- (1.501,0.237) -- (1.498,0.249) -- (1.494,0.262) -- (1.491,0.274) -- (1.487,0.287) -- (1.483,0.300) -- (1.479,0.312) -- (1.474,0.325) -- (1.469,0.338) -- (1.465,0.351) -- (1.459,0.364) -- (1.454,0.377) -- (1.449,0.391) -- (1.443,0.404) -- (1.437,0.417) -- (1.431,0.431) -- (1.425,0.444) -- (1.418,0.458) -- (1.411,0.471) -- (1.404,0.485) -- (1.397,0.498) -- (1.390,0.512) -- (1.382,0.526) -- (1.374,0.539) -- (1.366,0.553) -- (1.358,0.567) -- (1.350,0.580) -- (1.341,0.594) -- (1.332,0.608) -- (1.323,0.622) -- (1.314,0.636) -- (1.304,0.650) -- (1.295,0.664) -- (1.285,0.677) -- (1.275,0.691) -- (1.264,0.705) -- (1.254,0.719) -- (1.243,0.733) -- (1.232,0.747) -- (1.221,0.761) -- (1.210,0.775) -- (1.198,0.788) -- (1.187,0.802) -- (1.175,0.816) -- (1.163,0.830) -- (1.150,0.843) -- (1.138,0.857) -- (1.125,0.871) -- (1.113,0.884) -- (1.100,0.898) -- (1.086,0.911) -- (1.073,0.925) -- (1.059,0.938) -- (1.046,0.952) -- (1.032,0.965) -- (1.018,0.978) -- (1.003,0.991) -- (0.989,1.005) -- (0.974,1.018) -- (0.959,1.031) -- (0.945,1.044) -- (0.929,1.057) -- (0.914,1.069) -- (0.899,1.082) -- (0.883,1.095) -- (0.868,1.107) -- (0.852,1.120) -- (0.836,1.132) -- (0.820,1.144) -- (0.803,1.156) -- (0.787,1.169) -- (0.770,1.181) -- (0.754,1.192) -- (0.737,1.204) -- (0.720,1.216) -- (0.703,1.228) -- (0.686,1.239) -- (0.669,1.250) -- (0.651,1.262) -- (0.634,1.273) -- (0.616,1.284) -- (0.598,1.295) -- (0.581,1.305) -- (0.563,1.316) -- (0.545,1.326) -- (0.527,1.337) -- (0.509,1.347) -- (0.490,1.357) -- (0.472,1.367) -- (0.454,1.377) -- (0.435,1.387) -- (0.417,1.396) -- (0.398,1.406) -- (0.379,1.415) -- (0.361,1.424) -- (0.342,1.433) -- (0.323,1.442) -- (0.304,1.450) -- (0.286,1.459) -- (0.267,1.467) -- (0.248,1.476) -- (0.229,1.484) -- (0.210,1.491) -- (0.191,1.499) -- (0.172,1.507) -- (0.153,1.514) -- (0.133,1.521) -- (0.114,1.529) -- (0.095,1.535) -- (0.076,1.542) -- (0.057,1.549) -- (0.038,1.555) -- (0.019,1.561) -- (0.000,1.568) -- (-0.019,1.573) -- (-0.038,1.579) -- (-0.057,1.585) -- (-0.076,1.590) -- (-0.095,1.595) -- (-0.114,1.600) -- (-0.133,1.605) -- (-0.151,1.610) -- (-0.170,1.614) -- (-0.189,1.619) -- (-0.207,1.623) -- (-0.226,1.627) -- (-0.244,1.631) -- (-0.263,1.634) -- (-0.281,1.638) -- (-0.300,1.641) -- (-0.318,1.644) -- (-0.336,1.647) -- (-0.354,1.650) -- (-0.372,1.652) -- (-0.390,1.655) -- (-0.408,1.657) -- (-0.425,1.659) -- (-0.443,1.661) -- (-0.460,1.662) -- (-0.478,1.664) -- (-0.495,1.665) -- (-0.512,1.666) -- (-0.529,1.667) -- (-0.546,1.668) -- (-0.563,1.669) -- (-0.580,1.669) -- (-0.596,1.669) -- (-0.612,1.669) -- (-0.629,1.669) -- (-0.645,1.669) -- (-0.661,1.668) -- (-0.677,1.668) -- (-0.692,1.667) -- (-0.708,1.666) -- (-0.723,1.665) -- (-0.739,1.663) -- (-0.754,1.662) -- (-0.769,1.660) -- (-0.784,1.658) -- (-0.798,1.656) -- (-0.813,1.654) -- (-0.827,1.652) -- (-0.841,1.650) -- (-0.855,1.647) -- (-0.869,1.644) -- (-0.883,1.641) -- (-0.896,1.638) -- (-0.909,1.635) -- (-0.923,1.631) -- (-0.935,1.628) -- (-0.948,1.624);
\draw[scan] (-0.948,1.624) -- (-0.961,1.620) -- (-0.973,1.616) -- (-0.985,1.612) -- (-0.997,1.607) -- (-1.009,1.603) -- (-1.021,1.598) -- (-1.032,1.593) -- (-1.044,1.589) -- (-1.055,1.583) -- (-1.065,1.578) -- (-1.076,1.573) -- (-1.087,1.567) -- (-1.097,1.562) -- (-1.107,1.556) -- (-1.117,1.550) -- (-1.127,1.544) -- (-1.136,1.538) -- (-1.145,1.532) -- (-1.154,1.525) -- (-1.163,1.519) -- (-1.172,1.512) -- (-1.181,1.505) -- (-1.189,1.498) -- (-1.197,1.491) -- (-1.205,1.484) -- (-1.213,1.477) -- (-1.220,1.470) -- (-1.227,1.462) -- (-1.234,1.454) -- (-1.241,1.447) -- (-1.248,1.439) -- (-1.255,1.431) -- (-1.261,1.423) -- (-1.267,1.415) -- (-1.273,1.407) -- (-1.279,1.398) -- (-1.284,1.390) -- (-1.289,1.382) -- (-1.295,1.373) -- (-1.300,1.364) -- (-1.304,1.355) -- (-1.309,1.347) -- (-1.313,1.338) -- (-1.317,1.329) -- (-1.321,1.320) -- (-1.325,1.310) -- (-1.329,1.301) -- (-1.332,1.292) -- (-1.336,1.282) -- (-1.339,1.273) -- (-1.342,1.263) -- (-1.344,1.254) -- (-1.347,1.244) -- (-1.349,1.234) -- (-1.351,1.224) -- (-1.354,1.214) -- (-1.355,1.204) -- (-1.357,1.194) -- (-1.359,1.184) -- (-1.360,1.174) -- (-1.361,1.164) -- (-1.362,1.153) -- (-1.363,1.143) -- (-1.364,1.133) -- (-1.364,1.122) -- (-1.365,1.112) -- (-1.365,1.101) -- (-1.365,1.091) -- (-1.365,1.080) -- (-1.365,1.069) -- (-1.365,1.059) -- (-1.365,1.048) -- (-1.364,1.037) -- (-1.363,1.026) -- (-1.362,1.015) -- (-1.362,1.005) -- (-1.360,0.994) -- (-1.359,0.983) -- (-1.358,0.972) -- (-1.357,0.961) -- (-1.355,0.950) -- (-1.353,0.938) -- (-1.352,0.927) -- (-1.350,0.916) -- (-1.348,0.905) -- (-1.346,0.894) -- (-1.343,0.883) -- (-1.341,0.871) -- (-1.339,0.860) -- (-1.336,0.849) -- (-1.334,0.837) -- (-1.331,0.826) -- (-1.328,0.815) -- (-1.326,0.803) -- (-1.323,0.792) -- (-1.320,0.781) -- (-1.317,0.769) -- (-1.314,0.758) -- (-1.310,0.746) -- (-1.307,0.735) -- (-1.304,0.723) -- (-1.300,0.712) -- (-1.297,0.700) -- (-1.294,0.689) -- (-1.290,0.677) -- (-1.286,0.666) -- (-1.283,0.654) -- (-1.279,0.643) -- (-1.275,0.631) -- (-1.272,0.620) -- (-1.268,0.608) -- (-1.264,0.597) -- (-1.260,0.585) -- (-1.256,0.574) -- (-1.252,0.562) -- (-1.248,0.550) -- (-1.244,0.539) -- (-1.240,0.527) -- (-1.236,0.516) -- (-1.232,0.504) -- (-1.228,0.492) -- (-1.224,0.481) -- (-1.220,0.469) -- (-1.216,0.458) -- (-1.212,0.446) -- (-1.208,0.434) -- (-1.204,0.423) -- (-1.200,0.411) -- (-1.196,0.400) -- (-1.192,0.388) -- (-1.188,0.376) -- (-1.184,0.365) -- (-1.180,0.353) -- (-1.176,0.342) -- (-1.172,0.330) -- (-1.168,0.318) -- (-1.164,0.307) -- (-1.160,0.295) -- (-1.156,0.284) -- (-1.152,0.272) -- (-1.148,0.260) -- (-1.144,0.249) -- (-1.140,0.237) -- (-1.136,0.225) -- (-1.133,0.214) -- (-1.129,0.202) -- (-1.125,0.191) -- (-1.121,0.179) -- (-1.118,0.167) -- (-1.114,0.156) -- (-1.111,0.144) -- (-1.107,0.132) -- (-1.104,0.121) -- (-1.100,0.109) -- (-1.097,0.098) -- (-1.093,0.086) -- (-1.090,0.074) -- (-1.087,0.063) -- (-1.083,0.051) -- (-1.080,0.039) -- (-1.077,0.028) -- (-1.074,0.016) -- (-1.071,0.004) -- (-1.068,-0.007);
\draw[scanback] (-1.068,-0.007) -- (-1.065,-0.019) -- (-1.062,-0.031) -- (-1.059,-0.043) -- (-1.056,-0.054) -- (-1.053,-0.066) -- (-1.050,-0.078) -- (-1.048,-0.089) -- (-1.045,-0.101) -- (-1.042,-0.113) -- (-1.040,-0.125) -- (-1.037,-0.136) -- (-1.035,-0.148) -- (-1.032,-0.160) -- (-1.030,-0.172) -- (-1.028,-0.183) -- (-1.025,-0.195) -- (-1.023,-0.207) -- (-1.021,-0.219) -- (-1.018,-0.230) -- (-1.016,-0.242) -- (-1.014,-0.254) -- (-1.012,-0.266) -- (-1.010,-0.278) -- (-1.008,-0.290) -- (-1.006,-0.301) -- (-1.004,-0.313) -- (-1.002,-0.325) -- (-1.000,-0.337) -- (-0.998,-0.349) -- (-0.997,-0.361) -- (-0.995,-0.373) -- (-0.993,-0.385) -- (-0.991,-0.397) -- (-0.989,-0.409) -- (-0.988,-0.420) -- (-0.986,-0.432) -- (-0.984,-0.444) -- (-0.983,-0.456) -- (-0.981,-0.468) -- (-0.980,-0.480) -- (-0.978,-0.492) -- (-0.976,-0.504) -- (-0.975,-0.516) -- (-0.973,-0.528) -- (-0.972,-0.540) -- (-0.970,-0.552) -- (-0.969,-0.564) -- (-0.967,-0.576) -- (-0.965,-0.588) -- (-0.964,-0.600) -- (-0.962,-0.612) -- (-0.961,-0.624) -- (-0.959,-0.636) -- (-0.957,-0.648) -- (-0.956,-0.660) -- (-0.954,-0.672) -- (-0.953,-0.684) -- (-0.951,-0.696) -- (-0.949,-0.708) -- (-0.947,-0.720) -- (-0.946,-0.732) -- (-0.944,-0.744) -- (-0.942,-0.756) -- (-0.940,-0.768) -- (-0.938,-0.780) -- (-0.936,-0.792) -- (-0.934,-0.804) -- (-0.932,-0.816) -- (-0.930,-0.828) -- (-0.928,-0.840) -- (-0.926,-0.852) -- (-0.924,-0.864) -- (-0.922,-0.876) -- (-0.919,-0.888) -- (-0.917,-0.900) -- (-0.914,-0.912) -- (-0.912,-0.924) -- (-0.909,-0.935) -- (-0.907,-0.947) -- (-0.904,-0.959) -- (-0.901,-0.971) -- (-0.898,-0.983) -- (-0.895,-0.994) -- (-0.892,-1.006) -- (-0.889,-1.018) -- (-0.886,-1.029) -- (-0.883,-1.041) -- (-0.879,-1.052) -- (-0.876,-1.064) -- (-0.872,-1.075) -- (-0.869,-1.087) -- (-0.865,-1.098) -- (-0.861,-1.110) -- (-0.857,-1.121) -- (-0.853,-1.132) -- (-0.849,-1.144) -- (-0.844,-1.155) -- (-0.840,-1.166) -- (-0.836,-1.177) -- (-0.831,-1.188) -- (-0.826,-1.199) -- (-0.821,-1.210) -- (-0.816,-1.221) -- (-0.811,-1.232) -- (-0.806,-1.243) -- (-0.800,-1.254) -- (-0.795,-1.264) -- (-0.789,-1.275) -- (-0.784,-1.285) -- (-0.778,-1.296) -- (-0.772,-1.306) -- (-0.765,-1.317) -- (-0.759,-1.327) -- (-0.753,-1.337) -- (-0.746,-1.347) -- (-0.739,-1.357) -- (-0.732,-1.367) -- (-0.725,-1.377) -- (-0.718,-1.387) -- (-0.711,-1.397) -- (-0.703,-1.406) -- (-0.696,-1.416) -- (-0.688,-1.425) -- (-0.680,-1.435) -- (-0.672,-1.444) -- (-0.664,-1.453) -- (-0.655,-1.462) -- (-0.647,-1.471) -- (-0.638,-1.480) -- (-0.629,-1.489) -- (-0.620,-1.497) -- (-0.611,-1.506) -- (-0.601,-1.514) -- (-0.592,-1.522) -- (-0.582,-1.531) -- (-0.572,-1.539) -- (-0.562,-1.547) -- (-0.552,-1.554) -- (-0.542,-1.562) -- (-0.531,-1.570) -- (-0.520,-1.577) -- (-0.510,-1.585) -- (-0.499,-1.592) -- (-0.488,-1.599) -- (-0.476,-1.606) -- (-0.465,-1.613) -- (-0.453,-1.619) -- (-0.441,-1.626) -- (-0.429,-1.632) -- (-0.417,-1.638) -- (-0.405,-1.644) -- (-0.392,-1.650) -- (-0.380,-1.656) -- (-0.367,-1.662) -- (-0.354,-1.667) -- (-0.341,-1.673) -- (-0.328,-1.678) -- (-0.314,-1.683) -- (-0.301,-1.688) -- (-0.287,-1.693) -- (-0.273,-1.697) -- (-0.259,-1.701) -- (-0.245,-1.706) -- (-0.231,-1.710) -- (-0.217,-1.714) -- (-0.202,-1.717) -- (-0.187,-1.721) -- (-0.172,-1.724) -- (-0.157,-1.728) -- (-0.142,-1.731) -- (-0.127,-1.733) -- (-0.112,-1.736) -- (-0.096,-1.739) -- (-0.080,-1.741) -- (-0.064,-1.743) -- (-0.049,-1.745) -- (-0.032,-1.747) -- (-0.016,-1.748) -- (-0.000,-1.750) -- (0.016,-1.751) -- (0.033,-1.752) -- (0.050,-1.753) -- (0.066,-1.753) -- (0.083,-1.754) -- (0.100,-1.754) -- (0.117,-1.754) -- (0.134,-1.754) -- (0.152,-1.754) -- (0.169,-1.753) -- (0.187,-1.752) -- (0.204,-1.751) -- (0.222,-1.750);
\draw[scan] (0.222,-1.750) -- (0.240,-1.749) -- (0.257,-1.748) -- (0.275,-1.746) -- (0.293,-1.744) -- (0.311,-1.742) -- (0.329,-1.739) -- (0.348,-1.737) -- (0.366,-1.734) -- (0.384,-1.731) -- (0.402,-1.728) -- (0.421,-1.725) -- (0.439,-1.721) -- (0.458,-1.718) -- (0.476,-1.714) -- (0.495,-1.710) -- (0.514,-1.705) -- (0.532,-1.701) -- (0.551,-1.696) -- (0.570,-1.691) -- (0.588,-1.686) -- (0.607,-1.680) -- (0.626,-1.675) -- (0.645,-1.669) -- (0.664,-1.663) -- (0.682,-1.657) -- (0.701,-1.651) -- (0.720,-1.644) -- (0.739,-1.637) -- (0.757,-1.631) -- (0.776,-1.623) -- (0.795,-1.616) -- (0.814,-1.609) -- (0.832,-1.601) -- (0.851,-1.593) -- (0.870,-1.585) -- (0.888,-1.577) -- (0.907,-1.568) -- (0.926,-1.559) -- (0.944,-1.551) -- (0.962,-1.541) -- (0.981,-1.532) -- (0.999,-1.523) -- (1.017,-1.513) -- (1.036,-1.503) -- (1.054,-1.493) -- (1.072,-1.483) -- (1.090,-1.473) -- (1.108,-1.463) -- (1.126,-1.452) -- (1.143,-1.441) -- (1.161,-1.430) -- (1.178,-1.419) -- (1.196,-1.408) -- (1.213,-1.396) -- (1.230,-1.384) -- (1.248,-1.373) -- (1.265,-1.361) -- (1.282,-1.348) -- (1.298,-1.336) -- (1.315,-1.324) -- (1.332,-1.311) -- (1.348,-1.298) -- (1.364,-1.285) -- (1.380,-1.272) -- (1.396,-1.259) -- (1.412,-1.246) -- (1.428,-1.232) -- (1.444,-1.219) -- (1.459,-1.205) -- (1.474,-1.191) -- (1.489,-1.177) -- (1.504,-1.163) -- (1.519,-1.149) -- (1.534,-1.135) -- (1.548,-1.120) -- (1.562,-1.105) -- (1.577,-1.091) -- (1.590,-1.076) -- (1.604,-1.061) -- (1.618,-1.046) -- (1.631,-1.031) -- (1.644,-1.016) -- (1.657,-1.000) -- (1.670,-0.985) -- (1.683,-0.969) -- (1.695,-0.954) -- (1.707,-0.938) -- (1.719,-0.922) -- (1.731,-0.907) -- (1.743,-0.891) -- (1.754,-0.875) -- (1.766,-0.859) -- (1.776,-0.843) -- (1.787,-0.826) -- (1.798,-0.810) -- (1.808,-0.794) -- (1.818,-0.778) -- (1.828,-0.761) -- (1.838,-0.745) -- (1.847,-0.728) -- (1.857,-0.712) -- (1.866,-0.695) -- (1.874,-0.679) -- (1.883,-0.662) -- (1.891,-0.645) -- (1.899,-0.629) -- (1.907,-0.612) -- (1.915,-0.595) -- (1.922,-0.579) -- (1.929,-0.562) -- (1.936,-0.545) -- (1.943,-0.528) -- (1.949,-0.512) -- (1.956,-0.495) -- (1.961,-0.478) -- (1.967,-0.461) -- (1.973,-0.444) -- (1.978,-0.428) -- (1.983,-0.411) -- (1.988,-0.394) -- (1.992,-0.378) -- (1.997,-0.361) -- (2.001,-0.344) -- (2.004,-0.328) -- (2.008,-0.311) -- (2.011,-0.295) -- (2.014,-0.278) -- (2.017,-0.262) -- (2.020,-0.245) -- (2.022,-0.229) -- (2.024,-0.213) -- (2.026,-0.197) -- (2.028,-0.180) -- (2.030,-0.164) -- (2.031,-0.148) -- (2.032,-0.132) -- (2.032,-0.116) -- (2.033,-0.101) -- (2.033,-0.085) -- (2.033,-0.069) -- (2.033,-0.054) -- (2.033,-0.038) -- (2.032,-0.023) -- (2.031,-0.007) -- (2.030,0.008) -- (2.029,0.023) -- (2.027,0.038) -- (2.026,0.053) -- (2.024,0.068) -- (2.021,0.082) -- (2.019,0.097) -- (2.016,0.112) -- (2.014,0.126) -- (2.011,0.140) -- (2.007,0.154) -- (2.004,0.168) -- (2.000,0.182) -- (1.997,0.196) -- (1.993,0.210) -- (1.988,0.223) -- (1.984,0.237) -- (1.979,0.250) -- (1.975,0.263) -- (1.970,0.276) -- (1.964,0.289) -- (1.959,0.302) -- (1.954,0.314) -- (1.948,0.327) -- (1.942,0.339) -- (1.936,0.351) -- (1.930,0.363) -- (1.923,0.375) -- (1.917,0.387) -- (1.910,0.398) -- (1.903,0.410) -- (1.896,0.421) -- (1.889,0.432) -- (1.882,0.443) -- (1.874,0.454) -- (1.867,0.464) -- (1.859,0.475) -- (1.851,0.485) -- (1.843,0.495) -- (1.835,0.505) -- (1.826,0.515) -- (1.818,0.524) -- (1.809,0.534) -- (1.801,0.543) -- (1.792,0.552) -- (1.783,0.561) -- (1.774,0.570) -- (1.765,0.578) -- (1.756,0.587) -- (1.746,0.595) -- (1.737,0.603) -- (1.727,0.611) -- (1.718,0.619) -- (1.708,0.626) -- (1.698,0.634) -- (1.688,0.641) -- (1.678,0.648) -- (1.668,0.655) -- (1.658,0.661) -- (1.648,0.668) -- (1.638,0.674) -- (1.627,0.680) -- (1.617,0.686) -- (1.607,0.692) -- (1.596,0.697) -- (1.586,0.703) -- (1.575,0.708) -- (1.564,0.713) -- (1.554,0.718) -- (1.543,0.723) -- (1.532,0.727) -- (1.521,0.732) -- (1.510,0.736) -- (1.499,0.740) -- (1.488,0.744) -- (1.477,0.748) -- (1.466,0.751) -- (1.455,0.755) -- (1.444,0.758) -- (1.433,0.761) -- (1.422,0.764) -- (1.411,0.767) -- (1.400,0.769) -- (1.389,0.772) -- (1.378,0.774) -- (1.367,0.776) -- (1.356,0.778) -- (1.345,0.780) -- (1.334,0.781) -- (1.323,0.783) -- (1.312,0.784) -- (1.301,0.785) -- (1.290,0.786) -- (1.279,0.787) -- (1.268,0.788) -- (1.257,0.788) -- (1.246,0.789) -- (1.235,0.789) -- (1.224,0.789) -- (1.213,0.789) -- (1.202,0.789) -- (1.191,0.789) -- (1.181,0.789) -- (1.170,0.788) -- (1.159,0.788) -- (1.149,0.787) -- (1.138,0.786) -- (1.127,0.785) -- (1.117,0.784) -- (1.107,0.783) -- (1.096,0.781) -- (1.086,0.780) -- (1.075,0.778) -- (1.065,0.777) -- (1.055,0.775) -- (1.045,0.773) -- (1.035,0.771) -- (1.025,0.769) -- (1.015,0.767) -- (1.005,0.764) -- (0.995,0.762) -- (0.985,0.759) -- (0.975,0.757) -- (0.965,0.754) -- (0.956,0.751) -- (0.946,0.749) -- (0.937,0.746) -- (0.927,0.743) -- (0.918,0.740) -- (0.908,0.737) -- (0.899,0.733) -- (0.890,0.730) -- (0.881,0.727) -- (0.871,0.723) -- (0.862,0.720) -- (0.853,0.716) -- (0.845,0.713) -- (0.836,0.709) -- (0.827,0.705) -- (0.818,0.702) -- (0.809,0.698) -- (0.801,0.694) -- (0.792,0.690) -- (0.784,0.686) -- (0.775,0.682) -- (0.767,0.678) -- (0.759,0.674) -- (0.750,0.670) -- (0.742,0.666) -- (0.734,0.662) -- (0.726,0.658) -- (0.718,0.654) -- (0.710,0.650) -- (0.702,0.645) -- (0.694,0.641) -- (0.686,0.637) -- (0.678,0.633) -- (0.671,0.629) -- (0.663,0.624) -- (0.655,0.620) -- (0.648,0.616) -- (0.640,0.611) -- (0.632,0.607) -- (0.625,0.603) -- (0.618,0.599) -- (0.610,0.594) -- (0.603,0.590) -- (0.595,0.586) -- (0.588,0.582) -- (0.581,0.578) -- (0.574,0.573) -- (0.566,0.569) -- (0.559,0.565) -- (0.552,0.561) -- (0.545,0.557) -- (0.538,0.553) -- (0.530,0.549) -- (0.523,0.545) -- (0.516,0.541) -- (0.509,0.537) -- (0.502,0.533) -- (0.495,0.529) -- (0.488,0.525) -- (0.481,0.521) -- (0.474,0.518) -- (0.467,0.514) -- (0.460,0.510) -- (0.453,0.506) -- (0.446,0.503) -- (0.439,0.499) -- (0.432,0.496) -- (0.425,0.492) -- (0.417,0.489) -- (0.410,0.486) -- (0.403,0.482) -- (0.396,0.479) -- (0.389,0.476) -- (0.382,0.473) -- (0.374,0.470) -- (0.367,0.467) -- (0.360,0.464) -- (0.353,0.461) -- (0.345,0.458) -- (0.338,0.455) -- (0.330,0.452) -- (0.323,0.450) -- (0.315,0.447) -- (0.308,0.445) -- (0.300,0.442) -- (0.293,0.440) -- (0.285,0.437) -- (0.277,0.435) -- (0.269,0.433) -- (0.262,0.431) -- (0.254,0.429) -- (0.246,0.427) -- (0.238,0.425) -- (0.229,0.423) -- (0.221,0.421) -- (0.213,0.419) -- (0.205,0.418) -- (0.196,0.416) -- (0.188,0.415) -- (0.179,0.413) -- (0.171,0.412) -- (0.162,0.410) -- (0.153,0.409) -- (0.144,0.408) -- (0.135,0.407) -- (0.126,0.406) -- (0.117,0.405) -- (0.108,0.404) -- (0.098,0.403) -- (0.089,0.402) -- (0.080,0.401) -- (0.070,0.401) -- (0.060,0.400) -- (0.050,0.399) -- (0.041,0.399) -- (0.031,0.399) -- (0.021,0.398) -- (0.010,0.398) -- (-0.000,0.398) -- (-0.010,0.397) -- (-0.021,0.397);
\draw[scanback] (-0.021,0.397) -- (-0.032,0.397) -- (-0.042,0.397) -- (-0.053,0.397) -- (-0.064,0.397) -- (-0.075,0.398) -- (-0.086,0.398) -- (-0.098,0.398) -- (-0.109,0.398) -- (-0.120,0.399) -- (-0.132,0.399) -- (-0.144,0.400) -- (-0.156,0.400) -- (-0.168,0.401) -- (-0.180,0.401) -- (-0.192,0.402) -- (-0.204,0.403) -- (-0.217,0.403) -- (-0.229,0.404) -- (-0.242,0.405) -- (-0.254,0.406) -- (-0.267,0.407) -- (-0.280,0.408) -- (-0.293,0.409) -- (-0.307,0.410) -- (-0.320,0.411) -- (-0.333,0.412) -- (-0.347,0.413) -- (-0.361,0.414) -- (-0.374,0.415) -- (-0.388,0.416) -- (-0.402,0.418) -- (-0.416,0.419) -- (-0.431,0.420) -- (-0.445,0.421) -- (-0.459,0.423) -- (-0.474,0.424) -- (-0.489,0.425) -- (-0.503,0.426) -- (-0.518,0.428) -- (-0.533,0.429) -- (-0.548,0.430) -- (-0.563,0.432) -- (-0.579,0.433) -- (-0.594,0.435) -- (-0.609,0.436) -- (-0.625,0.437) -- (-0.640,0.439) -- (-0.656,0.440) -- (-0.672,0.442) -- (-0.688,0.443) -- (-0.704,0.444) -- (-0.720,0.446) -- (-0.736,0.447) -- (-0.752,0.448) -- (-0.768,0.450) -- (-0.785,0.451) -- (-0.801,0.452) -- (-0.818,0.454) -- (-0.834,0.455) -- (-0.851,0.456) -- (-0.867,0.457) -- (-0.884,0.459) -- (-0.901,0.460) -- (-0.918,0.461) -- (-0.935,0.462) -- (-0.952,0.463) -- (-0.969,0.464) -- (-0.986,0.465) -- (-1.003,0.466) -- (-1.020,0.467) -- (-1.037,0.468) -- (-1.054,0.469) -- (-1.071,0.470) -- (-1.089,0.471) -- (-1.106,0.472) -- (-1.123,0.473) -- (-1.140,0.473) -- (-1.158,0.474) -- (-1.175,0.475) -- (-1.192,0.475) -- (-1.210,0.476) -- (-1.227,0.476) -- (-1.244,0.477) -- (-1.262,0.477) -- (-1.279,0.478) -- (-1.296,0.478) -- (-1.314,0.478) -- (-1.331,0.478) -- (-1.348,0.479) -- (-1.366,0.479) -- (-1.383,0.479) -- (-1.400,0.479) -- (-1.417,0.479) -- (-1.434,0.479) -- (-1.451,0.478) -- (-1.468,0.478) -- (-1.485,0.478) -- (-1.502,0.478) -- (-1.519,0.477) -- (-1.536,0.477) -- (-1.553,0.476) -- (-1.569,0.476) -- (-1.586,0.475) -- (-1.603,0.474) -- (-1.619,0.474) -- (-1.635,0.473) -- (-1.652,0.472) -- (-1.668,0.471) -- (-1.684,0.470) -- (-1.700,0.469) -- (-1.716,0.467) -- (-1.732,0.466) -- (-1.748,0.465) -- (-1.763,0.464) -- (-1.779,0.462) -- (-1.794,0.461) -- (-1.809,0.459) -- (-1.824,0.457) -- (-1.840,0.456) -- (-1.854,0.454) -- (-1.869,0.452) -- (-1.884,0.450) -- (-1.898,0.448) -- (-1.913,0.446) -- (-1.927,0.444) -- (-1.941,0.442) -- (-1.955,0.440) -- (-1.969,0.437) -- (-1.982,0.435) -- (-1.996,0.433) -- (-2.009,0.430) -- (-2.022,0.427) -- (-2.035,0.425) -- (-2.048,0.422) -- (-2.060,0.419) -- (-2.073,0.417) -- (-2.085,0.414) -- (-2.097,0.411) -- (-2.109,0.408) -- (-2.120,0.405) -- (-2.132,0.401) -- (-2.143,0.398) -- (-2.154,0.395) -- (-2.165,0.392) -- (-2.176,0.388) -- (-2.186,0.385) -- (-2.196,0.381) -- (-2.206,0.378) -- (-2.216,0.374) -- (-2.226,0.371) -- (-2.235,0.367) -- (-2.244,0.363) -- (-2.253,0.359) -- (-2.262,0.355) -- (-2.271,0.352) -- (-2.279,0.348) -- (-2.287,0.344) -- (-2.295,0.339) -- (-2.302,0.335) -- (-2.310,0.331) -- (-2.317,0.327) -- (-2.324,0.323) -- (-2.330,0.318) -- (-2.337,0.314) -- (-2.343,0.310) -- (-2.349,0.305) -- (-2.354,0.301) -- (-2.360,0.296) -- (-2.365,0.292) -- (-2.370,0.287) -- (-2.374,0.283) -- (-2.379,0.278) -- (-2.383,0.273) -- (-2.387,0.269) -- (-2.390,0.264) -- (-2.394,0.259) -- (-2.397,0.254) -- (-2.400,0.250) -- (-2.402,0.245) -- (-2.405,0.240) -- (-2.407,0.235) -- (-2.409,0.230) -- (-2.410,0.225) -- (-2.412,0.220) -- (-2.413,0.215) -- (-2.413,0.210);
\draw[scan] (-2.413,0.210) -- (-2.414,0.205) -- (-2.414,0.200) -- (-2.414,0.195) -- (-2.414,0.190) -- (-2.414,0.185) -- (-2.413,0.180) -- (-2.412,0.175) -- (-2.411,0.170) -- (-2.409,0.165) -- (-2.407,0.160) -- (-2.405,0.155) -- (-2.403,0.150) -- (-2.400,0.145) -- (-2.398,0.140) -- (-2.395,0.135) -- (-2.391,0.129) -- (-2.388,0.124) -- (-2.384,0.119) -- (-2.380,0.114) -- (-2.376,0.109) -- (-2.371,0.104) -- (-2.366,0.099) -- (-2.361,0.094) -- (-2.356,0.089) -- (-2.350,0.084) -- (-2.345,0.079) -- (-2.339,0.074) -- (-2.333,0.069) -- (-2.326,0.064) -- (-2.319,0.059) -- (-2.313,0.054) -- (-2.305,0.050) -- (-2.298,0.045) -- (-2.290,0.040) -- (-2.283,0.035) -- (-2.275,0.030) -- (-2.266,0.025) -- (-2.258,0.021) -- (-2.249,0.016) -- (-2.240,0.011) -- (-2.231,0.007) -- (-2.222,0.002) -- (-2.213,-0.003) -- (-2.203,-0.007) -- (-2.193,-0.012) -- (-2.183,-0.016) -- (-2.173,-0.021) -- (-2.162,-0.025) -- (-2.152,-0.030) -- (-2.141,-0.034) -- (-2.130,-0.038) -- (-2.119,-0.043) -- (-2.107,-0.047) -- (-2.096,-0.051) -- (-2.084,-0.055) -- (-2.072,-0.059) -- (-2.060,-0.064) -- (-2.048,-0.068) -- (-2.036,-0.072) -- (-2.023,-0.076) -- (-2.011,-0.080) -- (-1.998,-0.084) -- (-1.985,-0.088) -- (-1.972,-0.091) -- (-1.959,-0.095) -- (-1.946,-0.099) -- (-1.932,-0.103) -- (-1.919,-0.106) -- (-1.905,-0.110) -- (-1.891,-0.114) -- (-1.877,-0.117) -- (-1.863,-0.121) -- (-1.849,-0.124) -- (-1.835,-0.128) -- (-1.820,-0.131) -- (-1.806,-0.134) -- (-1.791,-0.138) -- (-1.776,-0.141) -- (-1.762,-0.144) -- (-1.747,-0.147) -- (-1.732,-0.151) -- (-1.717,-0.154) -- (-1.702,-0.157) -- (-1.687,-0.160) -- (-1.671,-0.163) -- (-1.656,-0.166) -- (-1.641,-0.169) -- (-1.625,-0.172) -- (-1.610,-0.174) -- (-1.594,-0.177) -- (-1.579,-0.180) -- (-1.563,-0.183) -- (-1.547,-0.186) -- (-1.532,-0.188) -- (-1.516,-0.191) -- (-1.500,-0.193) -- (-1.484,-0.196) -- (-1.468,-0.199) -- (-1.453,-0.201) -- (-1.437,-0.204) -- (-1.421,-0.206) -- (-1.405,-0.208) -- (-1.389,-0.211) -- (-1.373,-0.213) -- (-1.357,-0.215) -- (-1.341,-0.218) -- (-1.325,-0.220) -- (-1.309,-0.222) -- (-1.293,-0.224) -- (-1.277,-0.227) -- (-1.261,-0.229) -- (-1.245,-0.231) -- (-1.230,-0.233) -- (-1.214,-0.235) -- (-1.198,-0.237) -- (-1.182,-0.239) -- (-1.166,-0.242) -- (-1.151,-0.244) -- (-1.135,-0.246) -- (-1.119,-0.248) -- (-1.103,-0.250) -- (-1.088,-0.252) -- (-1.072,-0.254) -- (-1.057,-0.256) -- (-1.041,-0.258) -- (-1.026,-0.260) -- (-1.011,-0.261) -- (-0.995,-0.263) -- (-0.980,-0.265) -- (-0.965,-0.267) -- (-0.950,-0.269) -- (-0.935,-0.271) -- (-0.920,-0.273) -- (-0.905,-0.275) -- (-0.890,-0.277) -- (-0.875,-0.279) -- (-0.860,-0.281) -- (-0.846,-0.283) -- (-0.831,-0.285) -- (-0.816,-0.287) -- (-0.802,-0.289) -- (-0.788,-0.291) -- (-0.773,-0.293) -- (-0.759,-0.295) -- (-0.745,-0.297) -- (-0.731,-0.299) -- (-0.717,-0.301) -- (-0.703,-0.303) -- (-0.690,-0.305) -- (-0.676,-0.307) -- (-0.662,-0.309) -- (-0.649,-0.312) -- (-0.635,-0.314) -- (-0.622,-0.316) -- (-0.609,-0.318) -- (-0.596,-0.320) -- (-0.583,-0.323) -- (-0.570,-0.325) -- (-0.557,-0.327) -- (-0.544,-0.330) -- (-0.531,-0.332) -- (-0.519,-0.335) -- (-0.506,-0.337) -- (-0.494,-0.340) -- (-0.482,-0.342) -- (-0.469,-0.345) -- (-0.457,-0.347) -- (-0.445,-0.350);
\draw[scanback] (-0.445,-0.350) -- (-0.433,-0.353) -- (-0.422,-0.355) -- (-0.410,-0.358) -- (-0.398,-0.361) -- (-0.387,-0.364) -- (-0.375,-0.367) -- (-0.364,-0.370) -- (-0.353,-0.373) -- (-0.341,-0.376) -- (-0.330,-0.379) -- (-0.319,-0.382) -- (-0.308,-0.385) -- (-0.298,-0.388) -- (-0.287,-0.391) -- (-0.276,-0.395) -- (-0.266,-0.398) -- (-0.255,-0.401) -- (-0.245,-0.405) -- (-0.235,-0.408) -- (-0.224,-0.412) -- (-0.214,-0.416) -- (-0.204,-0.419) -- (-0.194,-0.423) -- (-0.184,-0.427) -- (-0.174,-0.430) -- (-0.165,-0.434) -- (-0.155,-0.438) -- (-0.145,-0.442) -- (-0.136,-0.446) -- (-0.126,-0.450) -- (-0.117,-0.454) -- (-0.108,-0.459) -- (-0.098,-0.463) -- (-0.089,-0.467) -- (-0.080,-0.471) -- (-0.071,-0.476) -- (-0.062,-0.480) -- (-0.053,-0.485) -- (-0.044,-0.489) -- (-0.035,-0.494) -- (-0.026,-0.498) -- (-0.017,-0.503) -- (-0.009,-0.508) -- (0.000,-0.513) -- (0.009,-0.518) -- (0.017,-0.522) -- (0.026,-0.527) -- (0.034,-0.532) -- (0.043,-0.537) -- (0.051,-0.542) -- (0.060,-0.548) -- (0.068,-0.553) -- (0.077,-0.558) -- (0.085,-0.563) -- (0.093,-0.569) -- (0.102,-0.574) -- (0.110,-0.579) -- (0.118,-0.585) -- (0.126,-0.590) -- (0.135,-0.596) -- (0.143,-0.601) -- (0.151,-0.607) -- (0.159,-0.612) -- (0.168,-0.618) -- (0.176,-0.624) -- (0.184,-0.629) -- (0.192,-0.635) -- (0.201,-0.641) -- (0.209,-0.647) -- (0.217,-0.652) -- (0.225,-0.658) -- (0.234,-0.664) -- (0.242,-0.670) -- (0.250,-0.676) -- (0.258,-0.682) -- (0.267,-0.688) -- (0.275,-0.694) -- (0.284,-0.700) -- (0.292,-0.706) -- (0.300,-0.712) -- (0.309,-0.718) -- (0.317,-0.724) -- (0.326,-0.730) -- (0.334,-0.736) -- (0.343,-0.742) -- (0.352,-0.748) -- (0.360,-0.754) -- (0.369,-0.760) -- (0.378,-0.766) -- (0.387,-0.772) -- (0.396,-0.778) -- (0.405,-0.784) -- (0.414,-0.790) -- (0.423,-0.796) -- (0.432,-0.802) -- (0.441,-0.808) -- (0.450,-0.814) -- (0.459,-0.820) -- (0.469,-0.826) -- (0.478,-0.831) -- (0.487,-0.837) -- (0.497,-0.843) -- (0.506,-0.849) -- (0.516,-0.854) -- (0.526,-0.860) -- (0.535,-0.866) -- (0.545,-0.871) -- (0.555,-0.877) -- (0.565,-0.882) -- (0.575,-0.888) -- (0.585,-0.893) -- (0.595,-0.899) -- (0.605,-0.904) -- (0.616,-0.909) -- (0.626,-0.914) -- (0.637,-0.919) -- (0.647,-0.924) -- (0.658,-0.929) -- (0.668,-0.934) -- (0.679,-0.939) -- (0.690,-0.944) -- (0.701,-0.949) -- (0.712,-0.953) -- (0.723,-0.958) -- (0.734,-0.962) -- (0.745,-0.967) -- (0.757,-0.971) -- (0.768,-0.975) -- (0.779,-0.979) -- (0.791,-0.983) -- (0.802,-0.987) -- (0.814,-0.991) -- (0.826,-0.994) -- (0.838,-0.998) -- (0.849,-1.002) -- (0.861,-1.005) -- (0.873,-1.008) -- (0.885,-1.011) -- (0.898,-1.014) -- (0.910,-1.017) -- (0.922,-1.020) -- (0.934,-1.023) -- (0.947,-1.025) -- (0.959,-1.028) -- (0.972,-1.030) -- (0.985,-1.032) -- (0.997,-1.034) -- (1.010,-1.036) -- (1.023,-1.038) -- (1.036,-1.040) -- (1.048,-1.041) -- (1.061,-1.042) -- (1.074,-1.044) -- (1.088,-1.045) -- (1.101,-1.046) -- (1.114,-1.046) -- (1.127,-1.047) -- (1.140,-1.048) -- (1.154,-1.048) -- (1.167,-1.048) -- (1.180,-1.048) -- (1.194,-1.048) -- (1.207,-1.048) -- (1.221,-1.047) -- (1.234,-1.046) -- (1.248,-1.046) -- (1.261,-1.045) -- (1.275,-1.043) -- (1.289,-1.042) -- (1.302,-1.041) -- (1.316,-1.039) -- (1.329,-1.037) -- (1.343,-1.035) -- (1.357,-1.033) -- (1.371,-1.031) -- (1.384,-1.028) -- (1.398,-1.025) -- (1.412,-1.022) -- (1.425,-1.019) -- (1.439,-1.016) -- (1.453,-1.013) -- (1.466,-1.009) -- (1.480,-1.005) -- (1.494,-1.001) -- (1.507,-0.997) -- (1.521,-0.993) -- (1.535,-0.988) -- (1.548,-0.983) -- (1.562,-0.979) -- (1.575,-0.973) -- (1.588,-0.968) -- (1.602,-0.963) -- (1.615,-0.957) -- (1.629,-0.951) -- (1.642,-0.945) -- (1.655,-0.939) -- (1.668,-0.932) -- (1.681,-0.926) -- (1.694,-0.919) -- (1.707,-0.912) -- (1.720,-0.905) -- (1.733,-0.897) -- (1.746,-0.890) -- (1.758,-0.882) -- (1.771,-0.874) -- (1.783,-0.866) -- (1.796,-0.857) -- (1.808,-0.849) -- (1.820,-0.840) -- (1.832,-0.831) -- (1.844,-0.822) -- (1.856,-0.813) -- (1.868,-0.804) -- (1.880,-0.794) -- (1.891,-0.784) -- (1.903,-0.774) -- (1.914,-0.764) -- (1.925,-0.754) -- (1.937,-0.743) -- (1.947,-0.732) -- (1.958,-0.721) -- (1.969,-0.710) -- (1.980,-0.699) -- (1.990,-0.688) -- (2.000,-0.676) -- (2.010,-0.664) -- (2.020,-0.652) -- (2.030,-0.640) -- (2.040,-0.628) -- (2.049,-0.616) -- (2.059,-0.603) -- (2.068,-0.590) -- (2.077,-0.577) -- (2.086,-0.564) -- (2.094,-0.551) -- (2.103,-0.538) -- (2.111,-0.524) -- (2.119,-0.510) -- (2.127,-0.497) -- (2.135,-0.483) -- (2.142,-0.468) -- (2.150,-0.454) -- (2.157,-0.440) -- (2.164,-0.425);

  \end{scope}
  \node[lab, anchor=north] at (10.45,-4.85) {a path that never closes};
  \begin{axis}[snd, at={(12.85cm,-4.95cm)}, xmin=0, xmax=5.7, ymax=1.75, xlabel={frequency ÷ f₁ \quad (f₂ = 1.618 f₁)}]
    \addplot[ycomb, line width=0.8pt, ink3!70] table[x=freq, y=amp, restrict expr to domain={\thisrow{axis}}{1:1}] {wlm-lattice.dat};
    \addplot[ycomb, line width=1.4pt, ink] table[x=freq, y=amp, restrict expr to domain={\thisrow{axis}}{0:0}] {wlm-lattice.dat};
    \node[lab, anchor=north east, align=right] at (axis cs:5.7,1.78) {heavy: mixes · light: own};
  \end{axis}
\end{scope}

% =============== 02 · PUT THEM IN A ROW ===============
\def\y{-6.3}
\rowframe{\y}{02 \enspace PUT THEM IN A ROW}{a lattice in time: echoes at every j·T₁ + k·T₂}
\begin{scope}[shift={(0,\y)}]
  \node[lab] (in) at (0.1,-2.6) {in};
  \node[adder] (a1) at (0.75,-2.6) {+};
  \node[blk] (d1) at (1.95,-2.6) {delay T₁};
  \node[tap] (t1) at (2.95,-2.6) {};
  \node[adder] (a2) at (3.65,-2.6) {+};
  \node[blk] (d2) at (4.85,-2.6) {delay T₂};
  \node[tap] (t2) at (5.85,-2.6) {};
  \draw[wire,->] (in) -- (a1); \draw[wire,->] (a1) -- (d1); \draw[wire] (d1) -- (t1); \draw[wire,->] (t1) -- (a2);
  \draw[wire,->] (a2) -- (d2); \draw[wire] (d2) -- (t2); \draw[wire,->] (t2) -- (6.65,-2.6) node[lab, above left] {out};
  \draw[back] (t1) -- ++(0,-1.25) -| (a1);
  \draw[back] (t2) -- ++(0,-1.25) -| (a2);
  % the loop
  \node[lab] (lin) at (7.25,-3.2) {in};
  \draw[loop, spin] (8.45,-3.2) circle (0.62); \node[lab] at (8.45,-3.2) {T₁};
  \draw[loop, spin] (10.65,-3.2) circle (1.0);  \node[lab] at (10.65,-3.2) {T₂};
  \draw[wire,->] (lin) -- (7.83,-3.2); \draw[wire,->] (9.07,-3.2) -- (9.65,-3.2); \draw[wire,->] (11.65,-3.2) -- (12.25,-3.2);
  \node[lab, anchor=north] at (9.6,-4.85) {round one j times, the other k};
  \begin{axis}[snd, at={(12.85cm,-4.95cm)}, xmin=0, xmax=7.7, ymax=1.15, xlabel={time ÷ T₁ \quad (T₂ = 1.618 T₁)}]
    \addplot[ycomb, line width=1.4pt, ink] table[x=time, y=amp] {wlm-echoes.dat};
  \end{axis}
\end{scope}

% =============== 03 · BEND A LOOP ===============
\def\y{-12.6}
\rowframe{\y}{03 \enspace BEND A LOOP}{a stretched series: higher partials go round a shorter loop}
\begin{scope}[shift={(0,\y)}]
  \node[lab] (in) at (0.1,-2.6) {in};
  \node[adder] (a) at (0.85,-2.6) {+};
  \node[blk] (d) at (2.25,-2.6) {delay T};
  \node[blk] (ap) at (4.0,-2.6) {allpass};
  \node[tap] (t) at (5.15,-2.6) {};
  \draw[wire,->] (in) -- (a); \draw[wire,->] (a) -- (d); \draw[wire,->] (d) -- (ap); \draw[wire] (ap) -- (t);
  \draw[wire,->] (t) -- (6.65,-2.6) node[lab, above left] {out};
  \draw[back] (t) -- ++(0,-1.25) -| (a);
  \node[lab, anchor=north] at (3.2,-4.2) {the allpass delays high partials less};
  % the loop: each partial's own ring
  \begin{scope}[shift={(8.75,-3.2)}, scale=0.78]
    % panel 3: partial n goes round a loop of length T / sqrt(1 + B n^2), B exaggerated; generated by gen-loops.py
\draw[ring] (0,0) circle (1.700);
\draw[ink3, line width=0.3pt] (1.472,0.850) -- (1.819,1.050) node[lab, anchor=west] {n = 1};
\draw[ring] (0,0) circle (1.572);
\draw[ring] (0,0) circle (1.410);
\draw[ring] (0,0) circle (1.250);
\draw[ring] (0,0) circle (1.107);
\draw[ring] (0,0) circle (0.984);
\draw[ring] (0,0) circle (0.882);
\draw[ring] (0,0) circle (0.795);
\draw[ink3, line width=0.3pt] (0.689,-0.398) -- (1.819,-1.050) node[lab, anchor=west] {n = 8};

  \end{scope}
  \node[lab, anchor=north] at (9.6,-4.85) {ring n: T ÷ √(1 + B·n²)};
  \begin{axis}[snd, at={(12.85cm,-4.95cm)}, xmin=0, xmax=12.6, ymax=1.6, xlabel={frequency ÷ f₀ \quad (stiffness exaggerated)}]
    \foreach \k in {1,...,12} { \edef\t{\noexpand\draw[ink3!60, densely dashed, line width=0.4pt] (axis cs:\k,0) -- (axis cs:\k,1.1);} \t }
    \addplot[ycomb, line width=1.4pt, ink] table[x=freq, y=amp] {wlm-stretch.dat};
    \node[lab, anchor=north east, align=right] at (axis cs:12.6,1.62) {dashed: the harmonic grid};
  \end{axis}
\end{scope}

% =============== 04 · LET THEM PUSH BACK ===============
\def\y{-18.9}
\rowframe{\y}{04 \enspace LET THEM PUSH BACK}{locking with memory: two loops that keep themselves going}
\begin{scope}[shift={(0,\y)}]
  \node[adder] (aA) at (0.75,-1.9) {+};
  \node[blk] (dA) at (1.95,-1.9) {delay A};
  \node[blk] (nA) at (3.55,-1.9) {tanh};
  \node[tap] (tA) at (4.6,-1.9) {};
  \node[adder] (aB) at (0.75,-4.1) {+};
  \node[blk] (dB) at (1.95,-4.1) {delay B};
  \node[blk] (nB) at (3.55,-4.1) {tanh};
  \node[tap] (tB) at (4.6,-4.1) {};
  \draw[wire,->] (aA) -- (dA); \draw[wire,->] (dA) -- (nA); \draw[wire] (nA) -- (tA);
  \draw[wire,->] (aB) -- (dB); \draw[wire,->] (dB) -- (nB); \draw[wire] (nB) -- (tB);
  \draw[wire,->] (tA) -- (6.65,-1.9) node[lab, above left] {A out};
  \draw[wire,->] (tB) -- (6.65,-4.1) node[lab, above left] {B out};
  \draw[back] (tA) -- ++(0,0.75) -| (aA);
  \draw[back] (tB) -- ++(0,-0.75) -| (aB);
  \draw[wire,->] (tA) -- node[lab, pos=0.35, right] {c} ($(aB)+(0.6,0.5)$) -- (aB);
  \draw[wire,->] (tB) -- node[lab, pos=0.35, right] {c} ($(aA)+(0.6,-0.5)$) -- (aA);
  % the loop
  \draw[loop, spin] (8.0,-3.0) circle (0.62); \node[lab] at (8.0,-3.0) {220 Hz};
  \draw[loop, spin] (10.9,-3.0) circle (1.0);  \node[lab] at (10.9,-3.0) {136 Hz};
  \fill[ink] (8.62,-3.0) circle (1.6pt); \fill[ink] (9.9,-3.0) circle (1.6pt);
  \draw[ink, line width=0.7pt, decorate, decoration={coil, aspect=0.5, segment length=2mm, amplitude=1.5mm}] (8.62,-3.0) -- (9.9,-3.0);
  \node[lab] at (9.26,-2.5) {c};
  \node[lab, anchor=north] at (9.6,-4.85) {pull hard and both share one tone};
  \begin{axis}[snd, at={(13.35cm,-4.95cm)}, width=6.1cm, xmin=0, xmax=0.345, ymax=100, ytick={0,80}, ylabel={slips / s},
      xlabel={coupling c}, xtick={0,0.1,0.19,0.28}, xticklabels={0,0.1,0.19,0.28}]
    \addplot[const plot, line width=1.3pt, measure] table[x=c, y=slips] {wlm-up.dat};
    \addplot[const plot, line width=1.1pt, measure, densely dashed] table[x=c, y=slips] {wlm-down.dat};
    \node[font=\scriptsize, text=measure, anchor=west] at (axis cs:0.03,58) {up, solid};
    \node[font=\scriptsize, text=measure, anchor=west] at (axis cs:0.03,40) {down, dashed};
  \end{axis}
\end{scope}
\node[font=\scriptsize, anchor=north west, text=ink2, text width=19.6cm] at (0,-25.0) {Row 4, measured: going up the pair locks at c = 0.28; coming down it lets go at 0.19. Locked, both share 1077 Hz, near the 5th partial of one and the 8th of the other: 8 : 5, a Fibonacci ratio close to the golden one. Loops simplified: the model also has a lowpass and a DC blocker in each loop, drives the tanh at 1.5, and weights its own feedback by (1 − c). Rows 1–3 are computed, with stiffness exaggerated in row 3 so the stretch shows.};

% ---------- title block ----------
\draw[ink, line width=0.8pt] (-0.3,-26.1) rectangle (19.9,-27.0);
\draw[ink, line width=0.3pt] (9.6,-26.1) -- (9.6,-27.0) (14.6,-26.1) -- (14.6,-27.0);
\node[font=\scriptsize\bfseries, text=ink, anchor=west] at (-0.1,-26.55) {DSP IN LOOPING DIMENSIONS · WHERE THE LOOPS MEET};
\node[font=\scriptsize, text=ink, anchor=west] at (9.8,-26.55) {INK computed \quad \textcolor{measure}{RED measured}};
\node[font=\scriptsize, text=ink2, anchor=west] at (14.8,-26.55) {FIG · 2026-11-09 · PAPER};

\begin{scope}[on background layer]
  \fill[paper] ($(current bounding box.south west)+(-0.4,-0.4)$) rectangle ($(current bounding box.north east)+(0.4,0.4)$);
\end{scope}
\end{tikzpicture}
\end{document}
